\documentclass[11pt,a4paper]{article}
\usepackage[T1]{fontenc}
\usepackage[utf8]{inputenc}
\usepackage[french]{babel}
\frenchsetup{SmallCapsFigTabCaptions=false}
\usepackage{lmodern,microtype,amsmath}
\usepackage[margin=23mm]{geometry}
\usepackage{booktabs,array,tabularx,longtable}
\usepackage{xcolor,tikz,enumitem,float}
\usepackage{hyperref}
\usetikzlibrary{arrows.meta,positioning,fit,calc}
\definecolor{ink}{HTML}{1F242C}
\definecolor{muted}{HTML}{5C626C}
\definecolor{line}{HTML}{D9DEE5}
\definecolor{accent}{HTML}{AF0032}
\definecolor{rose}{HTML}{FDF0F4}
\hypersetup{colorlinks=true,linkcolor=black,citecolor=black,urlcolor=ink,
 pdftitle={LANCE et IoTChainBench : rapport d'ensemble du projet},
 pdfauthor={Projet LANCE -- rédaction par volets spécialisés, coordination Codex},pdflang={fr-FR}}
\urlstyle{same}
\def\UrlBreaks{\do\/\do\-\do\_\do\.\do\:}
\newcolumntype{Y}{>{\raggedright\arraybackslash}X}
\DeclareRobustCommand{\code}[1]{{\small\path{#1}}}
\newcommand{\revision}{2026-10-15 22:53 (Europe/Brussels, UTC+02:00)}
\newcommand{\source}[2]{\href{../../#1}{\nolinkurl{#1}}. #2}
\setlength{\parskip}{3pt}
\setlength{\emergencystretch}{2em}
\setlist{nosep,leftmargin=5mm,topsep=3pt}
\clubpenalty=10000 \widowpenalty=10000
\setcounter{tocdepth}{2}
\tikzset{
 box/.style={draw=line,fill=white,rounded corners=2pt,line width=.55pt,
 text=ink,align=center,inner sep=6pt,font=\sffamily\small},
 focal/.style={box,draw=accent,fill=rose,line width=.75pt},
 flow/.style={-{Latex[length=2mm,width=1.2mm]},draw=muted,line width=.65pt}
}

\title{LANCE et IoTChainBench\\[4pt]
\large Rapport d'ensemble du projet\\
Objectifs, conception, réalisations et état d'avancement}
\author{Projet LANCE\\\small Rédaction par volets spécialisés, coordination Codex}
\date{Dernière révision : \revision\\[3pt]
\small État documentaire au 15 octobre 2026 ; résultats historiques attribués}
\begin{document}
\maketitle
\begin{abstract}
LANCE est une plateforme de recherche pour l'audit de sécurité de réseaux
d'objets connectés autorisés, développée dans le contexte des travaux UMONS
du programme NATO SPS G6259. Le projet associe la représentation du réseau,
un moteur d'agents en six phases, une bibliothèque de contrôles, un laboratoire
automatisé, une application de suivi, le benchmark IoTChainBench et une chaîne
de spécialisation des modèles. Ce rapport présente ces ensembles, leurs
interfaces, leurs réalisations et leurs conditions de validation. Six rédactions
spécialisées examinent architecture, réseaux, sécurité, application, modèles et
évaluation à partir du code et des documents locaux. Les acquis comprennent
notamment un graphe NetworkX propre au run, quatre modes d'exécution,
des contrats de preuve, 29 scénarios publics, un comparateur de politiques
et un service multi-adaptateurs. Des audits historiques ont produit des scores
et des diagnostics ; les comparaisons contrôlées restent à achever et aucun
avantage du LLM ni gain attribuable à la spécialisation n'est établi par ce bilan.
Les qualifications logicielles et les travaux proposés gardent leur statut propre.
Les limites majeures concernent
la qualification des environnements, la preuve des transitions réseau,
l'indépendance des tests, l'attribution des gains aux modèles et l'achèvement
des campagnes contrôlées. La feuille de route relie les prochains livrables
à des critères de sortie vérifiables. Ce bilan ne vérifie pas en direct
le laboratoire et n'ajoute aucune nouvelle expérience d'audit ou d'entraînement.
\end{abstract}
\tableofcontents
\clearpage
\section{Contexte, finalité et trajectoire du projet}
\label{sec:contexte}
\subsection{Un projet d'audit IoT dans un programme de recherche plus large}
Le projet LANCE étudie l'automatisation d'audits de sécurité de réseaux d'objets
connectés. IoT signifie \emph{Internet of Things} ; OT, \emph{Operational
Technology}, désigne ici les composants et services liés aux systèmes
industriels. Une infrastructure de ville intelligente peut associer capteurs,
passerelles, équipements de réseau, courtiers de messages, services applicatifs
et systèmes de supervision. Sa sécurité dépend à la fois des propriétés de ces
composants et des relations qui les rendent accessibles.

Le document local préparé pour le jalon III du programme NATO SPS G6259,
\emph{Enhancing IoT Systems Security}, situe le travail de l'Université de Mons
dans les activités d'expérimentation virtuelle, d'évaluation automatisée et de
méthodologie de sécurité IoT/OT. Il décrit une plateforme Proxmox distincte du
testbed physique des partenaires, ainsi que le pipeline LANCE et le benchmark
IoTChainBench. Ce document daté du 1\textsuperscript{er} septembre 2026 donne le
contexte institutionnel ; il est utilisé comme déclaration de suivi, sans
supposer une validation indépendante de tous les résultats qu'il mentionne
\cite{intro-milestone}.

Le présent rapport développe le périmètre technique disponible dans le dépôt :
représentation du réseau et des connaissances, moteur d'agents, interfaces,
laboratoire, contrôles de sécurité, mesure expérimentale et spécialisation des
modèles. Les activités des autres partenaires fournissent le contexte du
programme. Leur avancement détaillé, leurs données financières et leurs
résultats propres relèvent de leurs documents de suivi.

Trois ensembles ont des fonctions différentes : LANCE conduit et restitue
l'audit ; IoTChainBench fournit des situations de test et leurs références
pour l'évaluer ; la chaîne d'entraînement prépare des variantes spécialisées
des modèles. Un audit peut employer un modèle déjà disponible auprès d'un
fournisseur, sans lancer d'entraînement. Le laboratoire sert à reconstruire
des situations connues ; sur un réseau existant, les observations restent
exploitables, mais l'absence de référence exhaustive empêche de mesurer toutes
les failles manquées \cite{intro-readme,arch-provider,sec-evaluation}.

\subsection{Le problème : relier une hypothèse à un effet observé}
Un audit utile doit dépasser l'énumération des ports ouverts et des versions
logicielles. Il doit déterminer quels services sont atteignables, quels droits
sont nécessaires, quels comportements sont réellement observés et quelles
conclusions restent indéterminées. Deux équipements portant le même logiciel
peuvent présenter des situations différentes selon leur configuration, leur
exposition réseau, l'identité utilisée ou l'état d'une session.

Les outils spécialisés savent effectuer des contrôles précis, mais leur
enchaînement et l'interprétation des résultats demandent un raisonnement sur le
contexte. Un modèle de langage de grande taille, ou LLM (\emph{Large Language
Model}), peut proposer des hypothèses et des actions à partir d'observations
hétérogènes. Il peut aussi produire une explication plausible sans preuve
suffisante, ignorer une erreur d'outil ou dépasser le budget avant de vérifier
une piste. Le projet vise donc à rendre ses décisions observables et
contestables, plutôt qu'à traiter sa réponse comme une vérité de sécurité
\cite{intro-readme,intro-protocol}.

LANCE signifie \emph{LLM Agent for Network Compromise Evaluation}. Le modèle
intervient dans une chaîne dont les autres éléments conservent des
responsabilités propres : les sondes observent, les contrats bornent la preuve,
le moteur orchestre et comptabilise, le benchmark confronte aux propriétés
attendues et le rapport rend les résultats consultables. Cette répartition est
centrale pour comprendre les réalisations du projet.

Une exécution, appelée \emph{run} dans le logiciel, regroupe une configuration
de départ, un périmètre autorisé, les actions réalisées et leurs fichiers de
résultat. L'agent est la boucle logicielle qui donne au modèle des observations,
reçoit ses propositions et fait exécuter les outils admis. Le modèle ne contacte
pas directement les équipements : le moteur traduit ses demandes en appels
contrôlés. Cette unité permet de relire une décision dans son contexte, de
distinguer deux essais et de conserver aussi les tentatives incomplètes
\cite{arch-pipeline,arch-runner}.

Concrètement, dans un scénario de laboratoire, l'opérateur choisit le réseau,
le modèle et les budgets, puis lance l'audit. Le moteur prépare l'environnement,
examine la topologie, recherche les services, analyse les faiblesses possibles,
vérifie les hypothèses et évalue les accès. Il conserve ensuite un rapport
lisible, les résultats structurés, les traces d'outils et la consommation.
L'opérateur peut examiner une confirmation, une tentative non concluante ou
une étape interrompue à partir de ces éléments. Ce parcours décrit le
fonctionnement prévu par le code ; il ne suppose ni découverte exhaustive
ni réussite de chaque phase \cite{arch-pipeline,arch-rapport,app-runs}.

\subsection{Objectifs complémentaires et critères de réussite}
Le projet possède plusieurs objectifs liés, mais évaluables séparément. Le
tableau~\ref{tab:intro-objectifs} les associe à un livrable et à une condition
de validation. Cette lecture évite qu'une interface aboutie, un modèle entraîné
ou un scénario défini soit utilisé comme preuve de toutes les autres dimensions.

\begin{table}[htbp]
\centering\small
\begin{tabularx}{\linewidth}{@{}>{\raggedright\arraybackslash}p{28mm}YY@{}}
\toprule
Objectif & Production attendue & Ce qui permet d'en apprécier la réussite\\
\midrule
Comprendre le réseau & Représentation des équipements, services, liens et chemins potentiels. & Provenance des informations et séparation entre topologie déclarée et connectivité observée.\\
Conduire un audit & Pipeline, outils, états et rapport. & Exécution traçable, résultats soutenus et échecs explicitement conservés.\\
Expérimenter & Laboratoire, scénarios, témoins et remises à zéro. & État initial vérifié et conditions reproductibles.\\
Mesurer & Vérités terrain, contrats probatoires et comparateurs. & Unités, dénominateurs, exclusions et versions connus.\\
Spécialiser les modèles & Données relues, adaptateurs, service et routage. & Gain attribuable à la spécialisation sur des évaluations appropriées.\\
Rendre le travail utilisable & API, tableau de bord, documentation et artefacts. & Parcours compréhensible et accès aux observations nécessaires à la décision.\\
\bottomrule
\end{tabularx}
\caption{Les objectifs couvrent un instrument logiciel, une infrastructure
expérimentale et une démarche scientifique. Un même livrable peut contribuer à
plusieurs objectifs sans suffire à les valider tous.}
\label{tab:intro-objectifs}
\end{table}

La question scientifique actuelle demande \textbf{dans quelles conditions et à
quel coût une politique avec LLM améliore l'audit face à une automatisation
conditionnelle compétente}. Elle porte sur les détections étayées, les faux
positifs, la fiabilité, l'adaptation et la consommation. Une question
complémentaire concerne la spécialisation d'un modèle local : un routage par
rôle et des adaptateurs entraînés améliorent-ils le comportement au-delà du
modèle de base et du seul guidage par prompts ? Ces questions appellent des
comparaisons distinctes et ne présument pas une réponse favorable
\cite{intro-protocol,intro-mission}.

Une politique est la règle qui choisit la prochaine action à partir des
observations disponibles. Une automatisation conditionnelle peut détecter un
port web, puis lancer les contrôles HTTP prévus pour ce service. Une politique
avec LLM peut proposer une autre vérification en interprétant la réponse reçue.
La comparaison examine si cette liberté améliore les résultats suffisamment
pour justifier ses appels au modèle, ses délais et ses erreurs. La présence
d'une explication en langage naturel ne constitue pas à elle seule le gain
recherché.

\subsection{Une évolution du graphe d'analyse vers la plateforme expérimentale}
Les documents historiques permettent de situer la trajectoire sans transformer
leurs affirmations anciennes en état actuel. Le rapport Q4 2025 décrit une
approche centrée sur Neo4j, un enrichissement par la base de vulnérabilités NVD
et une interaction en langage naturel. Le rapport Q1 2026 décrit ensuite le
passage à une infrastructure déclarative YAML, à un graphe NetworkX en mémoire
et à un pipeline d'agents. Le graphe courant est construit en mémoire avec
NetworkX et chaque exécution conserve son propre état ; le pipeline coordonne
les actions des phases et leurs productions. Les anciens serveurs et commandes
documentent cette trajectoire sans être supposés nécessaires au système actuel \cite{intro-q4,intro-q1}.

Les documents de mission placent également la spécialisation des modèles
locaux au centre d'un volet de travail : préparation des traces, adaptation
QLoRA, service multi-adaptateurs et routage par rôle. Le pipeline et le
laboratoire servent alors de cadre pour vérifier le comportement de ces
modèles. L'évolution récente de septembre et octobre 2026 renforce les contrats
de preuve, l'isolation des runs, l'analyse des échecs et les comparaisons avec
des politiques sans LLM \cite{intro-mission,intro-phases,intro-research}.

Cette trajectoire explique pourquoi le projet comprend plus qu'un agent
conversationnel : il associe un moteur, un modèle de réseau, une bibliothèque
de contrôles, une application de suivi, un banc d'essai et un outillage
d'apprentissage.

La figure~\ref{fig:intro-ensemble} relie ces ensembles avant leur description
détaillée. Elle distingue le moteur qui décide et coordonne, les moyens
d'observation, les productions consultables et les retours d'apprentissage.

\begin{figure}[H]
\centering
\begin{tikzpicture}
\node[box,text width=31mm,minimum height=16mm] (s) at (0,0)
 {Scénarios\\et topologies};
\node[focal,text width=34mm,minimum height=16mm] (p) at (5,0)
 {Moteur LANCE\\graphe et six phases};
\node[box,text width=31mm,minimum height=16mm] (m) at (10,0)
 {Modèles\\entraînement et service};
\node[box,text width=31mm,minimum height=16mm] (l) at (0,-2.3)
 {Laboratoire\\et contrôles réseau};
\node[box,text width=34mm,minimum height=16mm] (a) at (5,-2.3)
 {Artefacts\\observations et résultats};
\node[box,text width=31mm,minimum height=16mm] (e) at (10,-2.3)
 {Évaluation\\benchmark et diagnostics};
\node[box,text width=34mm,minimum height=16mm] (u) at (5,-4.6)
 {Application\\suivi et historique};
\draw[flow] (s) -- (p);
\draw[flow] (s) -- (l);
\draw[flow,<->] (p) -- (m);
\draw[flow,<->] (p.south west) -- node[above,sloped,font=\scriptsize]{outils} (l.north east);
\draw[flow] (p) -- (a);
\draw[flow] (a) -- (e);
\draw[flow] (a) -- (u);
\draw[flow] (e) -- node[right,font=\scriptsize,align=left]{retours de\\développement relus} (m);
\end{tikzpicture}
\caption{Vue d'ensemble du projet et des flux d'information. Les retours
vers l'apprentissage exigent une acceptation humaine et concernent les données
de développement autorisées. Les scénarios de test restent exclus de ce flux.}
\label{fig:intro-ensemble}
\end{figure}

\section{Méthode du rapport et lecture des états d'avancement}
\label{sec:methode-bilan}
\subsection{Six rédactions spécialisées et une intégration commune}
La révision répartit la rédaction entre six regards : architecture logicielle,
réseaux et infrastructure, sécurité et preuves, application et usages,
benchmark et méthode scientifique, modèles et entraînement. Chaque rédacteur
examine les sources de son domaine, produit une partie développée et identifie
ses références. L'intégration harmonise le vocabulaire, la progression et les
limites ; les explications nécessaires sont données au lieu de leur emploi. La relecture croisée et l'inspection du PDF vérifient ensuite la
cohérence du document dans son ensemble.

La copie locale est fondée sur le commit
\code{73369b03359d4facc8ec58a233c1623284e746d0}, avec des modifications et des
fichiers non versionnés. Le manifeste \code{sources-locales.json} conserve les
chemins et empreintes des sources utilisées. Les résultats sont attribués à
leurs comptes rendus, avec leurs dates. Aucune nouvelle campagne réseau,
connexion au laboratoire, génération par un fournisseur externe ni session
d'entraînement n'est lancée pour rédiger ce rapport.

\subsection{Ce que signifie un état annoncé}
Un composant \emph{implémenté} est présent dans le code et ses chemins
d'exécution inspectés. Une \emph{validation logicielle} exerce un comportement
dans des conditions définies, souvent avec outils ou réponses simulés. Une
\emph{observation de laboratoire} décrit un essai et son environnement. Une
\emph{proposition} décrit un travail à réaliser. Ces états apparaissent dans
les bilans des différentes parties.

Un contrat décrit des conditions explicites que le logiciel peut vérifier :
champs requis d'un résultat, observations admissibles pour une propriété,
règles de calcul ou critères de fin d'une phase. Il relie une déclaration
à des contrôles reproductibles. Annoncer une lecture MQTT anonyme demande,
par exemple, une observation de l'échange applicatif pertinent et de son
résultat. Le contrat reste limité aux cas qu'il reconnaît ; sa version est
conservée, car une évolution des règles peut changer l'interprétation du même
journal \cite{arch-runner,sec-preuves}.

Le rapport évalue la maturité par livrables et par preuves, sans calculer un
pourcentage global arbitraire. Un scénario défini n'est pas nécessairement
déployé ; un run achevé n'a pas nécessairement couvert toutes les propriétés ;
un score ancien ne devient pas comparable à un score nouveau après changement
de contrat. Les termes \emph{confirmé}, \emph{terminé} et \emph{validé} gardent
donc le périmètre qui leur donne un sens.

Les publications et anciens rapports donnent une provenance aux orientations
du projet. Ils ne remplacent pas la lecture du code courant et de ses
artefacts. Les limites d'accès ou de reproduction sont signalées ; une
inconnue reste une inconnue dans les tableaux d'avancement.


\section{Architecture logicielle et moteur d'agents}
\label{sec:architecture}

\subsection{Une orchestration aux responsabilités explicites}
LANCE organise l'audit autour d'un moteur qui possède le cycle de vie de
l'exécution. Il prépare les ressources, choisit les phases applicables,
vérifie leurs prérequis et conserve un état terminal. Le modèle intervient
dans les décisions qui lui sont déléguées ; il ne décide pas seul qu'un fichier
est valide, qu'un budget est disponible ou qu'une phase est terminée.
Cette séparation permet de modifier un modèle ou un prompt sans déplacer
automatiquement les responsabilités du contrôleur \cite{arch-pipeline,arch-runner}.

L'évolution du socle explique cette organisation. Le rapport du quatrième
trimestre 2025 décrivait Neo4j, MongoDB et une interrogation via MCP ; celui
du premier trimestre 2026 présentait le passage à YAML, NetworkX et ChromaDB.
Ce sont des descriptions historiques. Dans le code inspecté au 15 octobre
2026, le backend de graphe implémenté est NetworkX ; SQLite conserve notamment
le catalogue des fournisseurs et modèles, les runs et leur consommation.
L'abstraction du graphe autorise une autre implémentation, mais ne constitue
pas un adaptateur Neo4j opérationnel \cite{intro-q4,intro-q1,arch-graphe,arch-db}.

La figure~\ref{fig:arch-moteur} situe les responsabilités actuelles. Le pipeline
appelle des modules de phase explicitement enregistrés, sans découverte de
plugins. Un contexte de phase transmet les données et services nécessaires.
La boucle modèle--outils dialogue avec le fournisseur ; l'exécuteur applique
les contrôles communs et archive les observations. Le rapport et les phases
d'agrégation peuvent également travailler sans cette boucle lorsqu'une
transformation déterministe suffit. Le code conserve néanmoins des adaptateurs
et méthodes de compatibilité dans \code{Pipeline} : la séparation en modules
n'équivaut pas à six services indépendants \cite{arch-contexte,arch-dispatch}.

Concrètement, l'initialisation crée un dossier identifié, enregistre la
configuration et la révision du code, puis prépare la topologie ou le contexte
de découverte. Les phases retenues sont ordonnées ; sélectionner une phase
aval peut ajouter ses prérequis. Le moteur vérifie aussi la présence réelle
des livrables nécessaires : demander un rapport ne permet pas de contourner
une vérification absente. Après exécution ou interruption, un bloc de
finalisation tente de figer les consommations, effectuer le nettoyage prévu
et écrire les métadonnées terminales. L'échec d'une opération de finalisation
reste distinct de la cause initiale ; une conservation incomplète des mesures
peut dégrader le statut. Le résultat décrit ainsi le déroulement effectivement
enregistré, plutôt qu'une simple arrivée à la dernière instruction
\cite{arch-pipeline}.

\begin{figure}[htbp]
\centering
\begin{tikzpicture}
\node[focal,text width=29mm,minimum height=15mm] (p) at (0,0)
 {Pipeline\\cycle de vie};
\node[box,text width=31mm,minimum height=15mm] (ph) at (4.7,0)
 {Phases\\contextes explicites};
\node[box,text width=31mm,minimum height=15mm] (r) at (4.7,-1.8)
 {Boucle\\modèle--outils};
\node[box,text width=29mm,minimum height=15mm] (f) at (9.4,-1.8)
 {Adaptateur\\fournisseur};
\node[box,text width=31mm,minimum height=15mm] (e) at (4.7,-3.6)
 {Exécuteur\\d'outils};
\node[box,text width=29mm,minimum height=15mm] (s) at (0,-3.6)
 {État du run\\graphe et fichiers};
\draw[flow] (p) -- (ph);
\draw[flow] (p) -- (s);
\draw[flow] (ph) -- (r);
\draw[flow,<->] (r) -- (f);
\draw[flow,<->] (r) -- (e);
\draw[flow] (e) -- (s);
\end{tikzpicture}
\caption{Le pipeline possède le cycle de vie et l'état du run. Les phases
emploient les services du moteur ; les observations et fichiers restent
rattachés à cette exécution. Les flèches représentent des appels et échanges,
pas une séparation entre machines \cite{arch-contexte,arch-executeur}.}
\label{fig:arch-moteur}
\end{figure}

\subsection{Représenter le réseau et isoler l'état d'un run}
Le graphe distingue les équipements, leurs attributs et les relations de
connectivité. Le chargeur YAML construit des objets d'infrastructure puis un
\code{DiGraph} NetworkX. L'interface du backend propose consultation d'un
équipement, voisinage, surface exposée, chemins et export sérialisable.
Chaque opération possède sa sémantique : la recherche générique de chemins
utilise une vue non orientée, tandis que le calcul de risque et l'analyse
pondérée des attaques respectent la direction du graphe. Un chemin de
connectivité reste une propriété du modèle, sans prouver un passage exploité
\cite{arch-chargeur,arch-graphe,arch-risque,arch-chemins}.

Dans ce graphe, un nœud représente un équipement et un arc une relation de
connectivité déclarée. L'orientation distingue un passage de A vers B du passage
inverse ; elle importe lorsque routes et filtrage rendent les accès
asymétriques. La centralité d'intermédiarité souligne les équipements situés
sur de nombreux chemins calculés, par exemple une passerelle reliant plusieurs
zones. Elle aide à prioriser l'investigation. Une recherche générique qui ignore
l'orientation explore plutôt la structure des connexions ; son résultat ne
doit pas être interprété comme une autorisation de passage dans les deux sens.

Le vocabulaire d'enrichissement distingue quatre objets. NVD
(\emph{National Vulnerability Database}) est la base consultée ; CVE
(\emph{Common Vulnerabilities and Exposures}) identifie une vulnérabilité
publiée ; CPE (\emph{Common Platform Enumeration}) décrit un produit et sa
plateforme ; CVSS (\emph{Common Vulnerability Scoring System}) cote la sévérité.
Un identifiant CVE connu ne prouve pas son applicabilité à un équipement.

Le parcours d'infrastructure déclarée peut enrichir les équipements à partir
de correspondances CPE et des données NVD. Le score de risque, borné entre 0 et 10, combine
0,4 fois la composante vulnérabilités, 0,3 fois l'exposition et 0,3 fois la
centralité. L'exposition dépend des services et de la distance dirigée depuis
le nœud Internet ; la centralité repose sur l'intermédiarité dirigée. Ces poids
constituent la politique heuristique configurée \code{directed-v2}. L'analyse des chemins emploie des coûts dérivés de facteurs protocolaires
et d'exploitabilité ; une représentation pondérée permet aussi de comparer
des distances après modification de l'état. Ces calculs sont des heuristiques
d'aide à la priorisation. Ils ne mesurent ni une probabilité de compromission
calibrée ni un pivot observé \cite{arch-cve,arch-risque,arch-chemins,arch-pondere}.

Le comportement change selon l'entrée. Un scénario public fournit une
topologie et des services déclarés, sans précharger ses réponses d'évaluation :
les outils de risque et de chemins signalent alors que ces résultats ne sont
pas précalculés. Le mode découverte commence avec un contexte minimal et
l'enrichit par reconnaissance. Les outils exposent une topologie résumée,
des détails par équipement et une surface exploitable par les phases suivantes
\cite{arch-outils-graphe}.

Chaque \code{RunContext} possède son dossier et son \code{GraphContext}. Un
contexte d'exécution lie ce graphe aux appels, y compris dans les workers.
Les variables de phase sont copiées et présentées comme une correspondance
non modifiable ; des services explicites conservent les opérations autorisées.
Cette isolation évite de confondre deux runs, sans constituer un bac à sable
système. Les appels autonomes d'outils de graphe conservent une session de
compatibilité distincte ; le pipeline emploie son propre graphe
\cite{arch-contexte,arch-graphe-contexte}.

\subsection{Six phases avec des entrées et des sorties différentes}
\paragraph{Phase 1 : analyse de la représentation disponible.}
L'entrée comprend le périmètre et la topologie connue, éventuellement presque
vide en découverte. Le modèle consulte notamment la topologie, la surface
déclarée et les détails d'équipements pour hiérarchiser l'exploration. Il produit
\code{01_graph_analysis.md} ; une projection conserve les observations des
outils de graphe. La phase distingue les propriétés déclarées des observations
réseau qui restent à réaliser. Elle ne remplace donc pas la reconnaissance
\cite{arch-registre,arch-outils-graphe}.

\paragraph{Phase 2 : reconnaissance et couverture.}
À partir du périmètre et de l'analyse initiale, l'agent choisit des sondes ARP,
Nmap et les outils disponibles. Il recherche les hôtes, ports et versions,
puis confronte ces observations aux services attendus. Le journal alimente
\code{02_recon_evidence.json} et la restitution \code{02_recon.md}.
Le contrôleur vérifie la couverture requise avant sauvegarde. Une version
observée complète le graphe ; une sonde en échec reste signalée. La couverture
de ce contrat ne garantit pas la découverte de tous les services possibles
\cite{arch-recon,intro-phases}.

ARP recherche des équipements sur le segment local en reliant adresses IP
et adresses matérielles ; il ne découvre pas directement toutes les machines
derrière un routeur. Nmap interroge des destinations et ports pour identifier
les services joignables et, selon la sonde, des informations de version. La
couverture confronte les contrôles requis aux observations enregistrées : un
service SSH déclaré mais jamais interrogé reste un manque de couverture ; une
sonde exécutée en erreur reste un échec visible. Cette comptabilité distingue
le travail tenté des observations effectivement obtenues.

\paragraph{Phase 3 : observation puis analyse par équipement.}
Les scanners recueillent d'abord les résultats. Chaque analyste reçoit un
contexte d'appareil et peut employer des contrôles protocolaires bornés,
par exemple HTTP, TLS ou une recherche CVE ; aucun shell général n'est nécessaire
à ce parcours. Les résultats par appareil alimentent une agrégation déterministe
qui normalise, déduplique et conserve la provenance. Le livrable canonique
\code{03_vuln_analysis.json} et le registre brut restent distincts ; une
exclusion ne supprime pas la proposition originale. Les échecs des workers, c'est-à-dire des tâches d'analyse exécutées
séparément, restent visibles dans le statut de phase \cite{arch-analyse,arch-agregation}.

\paragraph{Phase 4 : contrôles ciblés des hypothèses.}
La file d'analyse détermine les tâches de vérification. Des workers sélectionnent
les outils pertinents, par exemple une requête HTTP ou un contrôle
d'identifiants, selon la cible et le contrat applicable. Les observations
attribuées aux tâches sont agrégées dans \code{04_exploitation.json}. Le
code reconstruit les statuts depuis les traces, indépendamment du verdict
écrit par le modèle. La découverte d'autres hôtes peut déclencher un suivi
d'analyse dans le parcours général ; le parcours comparatif borné conserve
son périmètre propre \cite{arch-verification,arch-runner}.

\paragraph{Phase 5 : actions et évaluation des accès.}
Le moteur prépare les points d'entrée, cibles et informations récupérées.
L'agent peut tester un accès, employer un identifiant ou réaliser une opération
autorisée sur un hôte. En profil complet, un marqueur de fin demande au
contrôleur de construire \code{05_intrusion.json} depuis le journal ; le
modèle ne réécrit pas toute la campagne. Une garde évite une clôture sans
tentative lorsque des entrées existent. Elle n'exige pas une réussite, et la
projection actuelle ne reconstruit pas les transitions réseau
\cite{arch-intrusion,intro-phases}.

Sur un même exemple SSH, la phase 4 cherche à confirmer un constat précis,
tel qu'un identifiant faible permettant une connexion. La phase 5 reçoit les
points d'entrée et informations récupérées, puis examine les accès et opérations
qu'ils rendent possibles dans le périmètre autorisé. Elle peut donc réemployer
un identifiant pour observer un accès effectif ; elle ne constitue pas une
seconde mesure de détection, et une connexion directe ne prouve pas un pivot.

\paragraph{Phase 6 : restitution à partir des faits.}
Les hypothèses, vérifications et observations d'accès deviennent des fiches.
Le modèle rédige leurs commentaires et recommandations sans outils, dans un
contexte borné. Le code assemble les faits et calcule la synthèse ; les
commentaires demeurent signalés comme non validés. Le rapport
\code{06_report.md} conserve les inconnues et les défauts de rédaction.
Les fiches déjà utilisables peuvent être réemployées si leurs sources, prompt,
modèle et réglages sont identiques. Cette reprise ne relance aucun contrôle
réseau \cite{arch-rapport,arch-rapport-code}.

Le livrable \code{06_report.md} est un document Markdown consultable dans
l'application. Il réunit une synthèse, des tableaux factuels, les fiches des
hypothèses, les observations d'accès et les limites. Une fiche sans test ou
non concluante conserve cet état ; les fichiers structurés et le journal
\code{tool_calls.jsonl} permettent d'en examiner les sources.

\subsection{Boucle modèle--outils, transactions et profils}
Une phase construit son prompt, la consigne transmise au modèle, à partir
de modèles de texte, d'inclusions
communes et des variables du run. Le catalogue décrit les outils par nom,
description et schéma d'arguments ; les définitions YAML peuvent produire des
adaptateurs de sous-processus ou référencer un gestionnaire Python. Les
définitions matérielles fournissent des propositions de commandes, sans prouver
qu'un équipement physique a été piloté. Les outils indisponibles sont filtrés
avant présentation au modèle \cite{arch-prompts,arch-outils}.

La boucle transmet l'historique au fournisseur, reçoit une réponse ou des
appels structurés, exécute les fonctions sélectionnées et réinjecte leurs
résultats. L'exécuteur commun attribue un identifiant, une cible, une durée
et une origine aux observations. Une sauvegarde est transactionnelle : la
soumission est archivée, validée puis promue ; un rejet revient au modèle
comme erreur d'outil. Un texte final sans sauvegarde acceptée ne suffit pas.
Les limites de tours et les gardes contre répétitions ou réponses sans action
bornent cette conversation \cite{arch-provider,arch-executeur,arch-runner}.

Pour un service web, le modèle peut demander une requête vers une adresse,
un port et un chemin précis au moyen d'arguments structurés. L'exécuteur
vérifie ces arguments, appelle le client HTTP et conserve la réponse. Le
modèle reçoit alors une observation sur laquelle appuyer la décision suivante.
S'il propose un constat, sa sauvegarde doit satisfaire le schéma et les règles
applicables. La demande d'outil, sa réponse et le constat qui l'interprète
sont trois objets distincts. Ce découpage aide à distinguer une erreur de
décision, une erreur d'outil et un refus du service audité.

Le tableau~\ref{tab:arch-profils} précise trois choix souvent confondus.
\code{auto} résout un profil ; il n'ajoute pas un troisième moteur.
La priorité va au choix explicite du run, puis au réglage du modèle, puis
aux paramètres actifs ou totaux connus. Le seuil courant est de 32 milliards :
au plus ce seuil, le profil est compact ; sans métadonnées, il est complet.
Cette règle d'ingénierie n'est pas une mesure de capacité du modèle
\cite{arch-profils}. Par exemple, les plafonds courants de phase 1 sont
20 tours et 4\,096 jetons de sortie par réponse en complet, contre
12 tours et 2\,048 jetons en compact ; ils ne décrivent pas la fenêtre
de contexte effectivement acceptée par le serveur.

\begin{table}[htbp]
\centering\small
\begin{tabularx}{\linewidth}{@{}p{20mm}YY@{}}
\toprule
Choix & Effet sur l'exécution & Limite d'interprétation\\
\midrule
\code{full} & Surface autorisée par phase, budgets complets ; choix d'outils plus étendu. & Aucun privilège supplémentaire au-delà du périmètre du run.\\
\code{compact} & Outils routés et budgets adaptés ; observations présentées sous forme réduite. & La réduction peut retirer du contexte utile ; les sources annexes restent nécessaires.\\
\code{auto} & Résolution traçable depuis le réglage et les métadonnées du modèle. & Le nombre de paramètres ne garantit ni compétence ni fiabilité.\\
\bottomrule
\end{tabularx}
\caption{Les profils modifient la présentation et les ressources, sans changer
la vérité terrain ni l'évaluateur \cite{arch-profils}.}
\label{tab:arch-profils}
\end{table}

\subsection{Maîtriser le contexte et reprendre les productions incomplètes}
L'enjeu n'est pas seulement la longueur d'un prompt initial. Une conversation
réexpédie aussi les résultats précédents ; multiplier équipements et sondes
augmente le texte traité et le risque de restitution incomplète. Le code
réduit ce couplage par des contextes spécialisés. La phase 3 sépare scan,
analyse d'appareil et agrégation. La reconnaissance peut, pour une topologie
déclarée éligible, séparer un balayage initial et des lots de couverture,
puis les fusionner depuis le journal. Le mode aveugle commence sans topologie
fournie ; le mode scellé ajoute une séparation technique entre l'agent et les
réponses de l'évaluateur. Ces modes conservent le parcours simple \cite{arch-analyse,arch-recon-lots}.

Plusieurs récupérations répondent à des signatures précises. Une phase 1
complète tronquée peut reprendre seulement la rédaction depuis les faits
de graphe enregistrés. Pour une analyse d'appareil tronquée, la phase 3 peut
répartir la restitution en blocs de services, sauvegardés séparément avant
assemblage. Le contrôleur refuse la promotion si un bloc requis manque.
La phase 6 retente une fiche tronquée, puis conserve explicitement son
incomplétude. Les reprises restent bornées ; elles ne certifient pas la
capacité effective du serveur distant \cite{arch-reprise-graphe,arch-reprise-blocs,arch-rapport}.

\paragraph{Exemple de fonctionnement, sans résultat expérimental.}
Considérons un équipement exposant SSH, HTTPS et MQTT. Les sondes peuvent
avoir abouti alors que la réponse d'analyse s'interrompt avant sauvegarde.
La récupération par blocs permet de traiter ces observations dans des
productions plus petites, sans réexécuter les sondes. Si la restitution MQTT
reste invalide, l'appareil demeure incomplet ; deux blocs acceptés ne deviennent
pas automatiquement une analyse complète. Le bénéfice attendu est de
localiser l'échec et de conserver le travail réutilisable, pas de garantir
moins de coût ou de faux négatifs \cite{arch-reprise-blocs}.

\subsection{Connaissances, fournisseurs, ressources et état de ce volet}
La mémoire comprend plusieurs objets différents. Les fichiers du run assurent
la continuité entre phases. Les fiches méthodologiques Markdown peuvent être
chargées selon leurs étiquettes. ChromaDB fournit une recherche sémantique
sur connaissances et historiques. Un embedding représente un texte par un
vecteur numérique afin de retrouver des contenus proches ; sa production
utilise ici le client Voyage. L'emploi de cette recherche dépend des services et identifiants configurés. Une similarité textuelle
n'établit pas qu'une CVE concerne la version observée : une classification CPE
séparée conserve les cas compatibles, conditionnels, indéterminés et
incompatibles \cite{arch-connaissances,arch-cve}.
Le stockage vectoriel prévoit un repli éphémère si son ouverture persistante
échoue ; disposer de l'API de recherche ne garantit donc pas, à lui seul,
la conservation durable d'une mémoire \cite{arch-stockage}.

Ces mécanismes changent dans le benchmark. La recherche CVE utilise un
instantané versionné, sans consulter le cache mutable ni interroger NVD en
cas d'absence. La recherche d'historique et les collections persistantes hors
méthodologie sont refusées ; les résultats du benchmark ne sont pas réinjectés
en mémoire épisodique. Hors de ce cadre, le parcours général conserve une
recherche en cache puis une interrogation distante. La reproductibilité dépend
donc aussi du mode de connaissances, pas seulement du nom du modèle
\cite{arch-connaissances,arch-memoire,arch-executeur}.

L'interface fournisseur unifie des endpoints compatibles avec le format
OpenAI, configurés notamment dans SQLite. Un modèle peut être choisi par
phase. Le contrôle d'admission sérialise, lorsque configuré, les requêtes
partageant une origine ; les outils ne gardent pas ce verrou. Cela aide un
serveur local contraint, sans créer de capacité GPU. Le transport distingue
les reprises d'une même requête des continuations du raisonnement : une reprise
réseau ne rejoue pas elle-même les outils \cite{arch-provider,arch-transport,arch-admission}.

Le suivi agrège jetons, appels, erreurs, temps des agents et temps écoulé.
Un jeton est une unité de texte traitée par le modèle : mot, fragment de mot
ou symbole. Des observations réexpédiées dans plusieurs requêtes peuvent donc
être comptabilisées plusieurs fois.
Les tarifs retenus sont conservés par phase avec leur provenance et leur
caractère estimé. La limite financière porte sur la consommation observée :
des appels déjà engagés peuvent encore aboutir. Elle ne réserve pas à l'avance
toute facture future et ne représente pas le coût matériel d'un modèle local.
La finalisation tente de préserver consommation, cause d'arrêt et diagnostics,
même lorsqu'une étape échoue \cite{arch-cout,arch-pipeline}.

Les acquis de ce volet sont un moteur commun, des contextes propres au run,
des sorties structurées et plusieurs reprises ciblées. Ses limites restent
la dépendance au comportement des modèles, à la disponibilité des outils et
aux ressources de calcul ; le découpage ne démontre pas une scalabilité
générale. Les prochaines étapes proposées sont de stabiliser les contrats
de contexte, qualifier leur couverture et comparer la fiabilité des profils
sur des charges définies. Elles devront préserver les statuts incomplets
et les coûts des tentatives, plutôt que compter seulement les rapports produits.


\section{Réseaux, laboratoire et infrastructure d'exécution}
\label{sec:net-infrastructure}

\subsection{Un laboratoire reproductible au service de l'audit}

Le volet réseau de LANCE fournit l'environnement dans lequel les décisions de
l'agent deviennent des observations vérifiables. Il doit rendre possible une
question concrète : après avoir réinstallé un scénario donné, le même service,
la même exposition et les mêmes restrictions d'accès sont-ils présents ? Sans
cette maîtrise, une absence de résultat pourrait provenir du laboratoire plutôt
que de la méthode d'audit. L'infrastructure constitue donc une partie de
l'instrument scientifique, avec ses propres contrats et limites.

La cible retenue est le laboratoire \textbf{nato / pve-nato}. Proxmox héberge
les charges expérimentales et le contrôleur nommé \code{nato-master}. Les
playbooks distinguent les conteneurs LXC Debian, utilisés pour les services,
et une machine virtuelle OpenWrt, utilisée comme routeur de scénario. Des
modèles préalablement préparés servent de base au clonage. Cette approche
réduit le coût de reconstruction et permet de combiner plusieurs rôles sans
disposer d'un équipement physique par rôle. Le contrôleur lui-même est
provisionné comme conteneur LXC ; le terme historique de « VM maître » ne
désigne donc pas une seconde machine physique \cite{net-master,net-deploy}.

Proxmox fournit la couche de virtualisation qui héberge les charges du
scénario. Un conteneur LXC isole un environnement Linux en partageant le noyau
de l'hôte ; une machine virtuelle possède son propre système, comme pour le
rôle OpenWrt. Ansible applique les procédures de création et de configuration
décrites dans des playbooks. Le bridge joue le rôle de commutateur virtuel
entre les interfaces. Cette chaîne permet de décrire un laboratoire dans des
fichiers puis de le reconstruire ; sa répétabilité dépend des modèles système,
des versions de paquets et des contrôles d'état initial.

Trois responsabilités doivent être distinguées. Le \emph{plan de commande}
correspond à l'opérateur, à l'API et aux workflows qui demandent une opération.
Le \emph{contrôleur} exécute LANCE, appelle Ansible et dialogue avec Proxmox pour
créer, configurer et retirer les charges. Le \emph{réseau d'essai} contient les
services effectivement interrogés par les outils. Le contrôleur possède une
interface de gestion et une interface vers le bridge de benchmark. L'accès
distant documenté utilise Tailscale ; les règles de sortie installées séparent
les besoins de gestion des échanges d'audit. Cette séparation décrite par le
provisionnement doit être vérifiée sur la cible avant une campagne : une
configuration dans Git ne démontre pas son application effective
\cite{net-master,net-access}.

La figure~\ref{fig:net-laboratoire} résume ces responsabilités. Le laboratoire
personnel nommé \code{pve} est explicitement exclu du projet. Une panne de
nato ne justifie aucun déplacement de scénario, de secret ou de configuration
vers cette autre infrastructure. L'identité de la cible prime sur la présence
d'un inventaire ou sur la possibilité technique d'établir une connexion
\cite{net-scope}.

\begin{figure}[H]
\centering
\begin{tikzpicture}
\node[box,text width=31mm,minimum height=15mm] (op) at (0,0)
 {Opérateur\\Interface et CLI};
\node[focal,text width=34mm,minimum height=15mm] (master) at (5,0)
 {nato-master\\LANCE et Ansible};
\node[box,text width=31mm,minimum height=15mm] (pve) at (10,0)
 {pve-nato\\Proxmox};
\node[box,text width=43mm,minimum height=15mm] (lab) at (5,-2.1)
 {Réseau d'essai partagé\\Bridge et VLAN selon scénario};
\node[box,text width=31mm,minimum height=15mm] (iot) at (0,-4.2)
 {Services IoT / IT\\MQTT, web, SSH};
\node[box,text width=31mm,minimum height=15mm] (ot) at (5,-4.2)
 {Cellule OT\\Services et simulateurs};
\node[box,text width=31mm,minimum height=15mm] (seg) at (10,-4.2)
 {Zones et relais\\Objectifs multi-hop};
\draw[flow] (op) -- node[above,font=\scriptsize]{demande} (master);
\draw[flow] (master) -- node[above,font=\scriptsize]{API} (pve);
\draw[flow] (master) -- node[right,font=\scriptsize]{audite} (lab);
\draw[flow] (pve) -- node[right,font=\scriptsize]{héberge} (lab);
\draw[flow] (lab) -- (iot);
\draw[flow] (lab) -- (ot);
\draw[flow] (lab) -- (seg);
\end{tikzpicture}
\caption{Organisation logique de l'infrastructure implémentée. Les trois
branches inférieures représentent des compositions possibles ; elles ne
décrivent pas trois laboratoires disponibles simultanément. Le dessin ne
représente ni un état déployé vérifié au 15 octobre 2026, ni les règles exactes
de connectivité de chaque scénario.}
\label{fig:net-laboratoire}
\end{figure}

\subsection{Topologies, zones et signification d'un pivot}

Le corpus couvre plusieurs organisations de réseau. Une topologie plate place
les services dans une surface directement explorable ; une étoile organise les
dépendances autour d'une passerelle ; les configurations de bâtiment, de ville
connectée ou d'edge/cloud distribuent des rôles fonctionnels. Les configurations
IT/IoT/OT introduisent des zones de confiance différentes. Leur intérêt est de
faire varier l'information disponible, la diversité des services et les chemins
nécessaires pour atteindre une ressource. Le nom d'une zone décrit une fonction
expérimentale ; il ne suffit pas à établir un filtrage effectif
\cite{net-scenarios,net-attacks}.

La segmentation doit être lue à plusieurs niveaux. Un VLAN sépare des domaines
de communication sur le bridge ; une route indique comment joindre un autre
sous-réseau ; un pare-feu peut autoriser ou interdire ce passage. Un hôte à
plusieurs interfaces peut relier deux segments sans que tous les services soient
accessibles depuis le point de départ. Le playbook de déploiement configure
ces propriétés au moyen d'interfaces étiquetées, d'adresses secondaires et,
pour certains rôles, de l'absence de passerelle. S13 possède une configuration
particulière de bridge conscient des VLAN et de routes vers les zones. Les
scénarios de pivots récents emploient aussi des interfaces temporaires de
préparation : le contrôle de leur retrait et de l'absence d'accès direct relève
de l'injection et de la vérification \cite{net-deploy,net-verify}.

S20 fournit un exemple plus précis qu'une simple étiquette « réseau segmenté ».
Sa topologie définit un point d'entrée, un relais à plusieurs interfaces et une
ressource finale dans un second segment isolé. Les liens attendus imposent un
accès SSH vers le premier hôte, puis un accès SSH vers le relais, enfin un accès
HTTP vers la ressource. Les variantes S24--S27 définissent des branches, des
chemins asymétriques ou une cascade plus profonde. Ce sont des objectifs du
support expérimental, dont les playbooks comportent des contrôles de chemin
direct ; leur présence ne démontre pas que LANCE les a parcourus
\cite{net-topology,net-verify}.

Un \emph{pivot réseau} exige une action depuis un accès précédemment obtenu
vers une cible devenue joignable grâce à cet accès. Récupérer un secret sur un
site puis l'employer sur un autre service directement accessible constitue une
dépendance logique, sans prouver un changement de position réseau. De même,
plusieurs opérations d'API ne créent pas plusieurs sauts. Cette distinction
est décisive pour le projet : le pipeline courant ne produit pas la provenance
causale nécessaire pour mesurer ses transitions réseau. Les noms « pivot » ou
« multi-hop » ne doivent donc pas être convertis en résultats de compromission
\cite{net-attacks}.

\subsection{Familles de services et degré de réalisme}

La messagerie MQTT suit un modèle de publication et de souscription : un
client publie dans un canal nommé \emph{topic} et d'autres clients s'y abonnent.
Les ACL sont les règles qui autorisent ces opérations. TLS protège le
transport ; mTLS ajoute une authentification mutuelle par certificats.
OTA désigne une mise à jour distante. Ces propriétés sont séparées : chiffrer
un échange ne garantit pas que le client possède le droit de lire ou de modifier
la ressource.


Le laboratoire représente l'IoT par ses dépendances réseau et applicatives :
messagerie entre équipements, interfaces d'administration, bases de données,
passerelles, vidéo et systèmes industriels. Certains composants sont des
services courants configurés pour le benchmark ; d'autres sont des simulateurs
déterministes conçus pour isoler une propriété. Cette diversité est utile, mais
elle impose de décrire la fidélité de chaque expérience. Le
tableau~\ref{tab:net-familles} donne les principales familles sans leur attribuer
une couverture expérimentale globale \cite{net-attacks,net-verify}.

\begin{table}[H]
\centering\small
\begin{tabularx}{\linewidth}{@{}>{\raggedright\arraybackslash}p{26mm}YYY@{}}
\toprule
Famille & Services représentés & Contrôles pertinents & Limites du support\\
\midrule
Messagerie IoT & MQTT/Mosquitto ; CoAP dans certaines compositions & Connexion, lecture d'un sujet, publication et authentification, vérifiées séparément & Une opération admise ne décrit pas toutes les ACL ; aucun réseau radio réel n'en découle.\\
Administration et données & SSH, HTTP, FTP, SNMP, MariaDB, Redis, Node-RED & Authentification, accès à une ressource, exposition de configuration et refus attendu & Les configurations de laboratoire ne représentent pas toutes les versions ou installations d'un produit.\\
Vidéo et bâtiment & Façades caméra HTTP, rôle NVR, interfaces de bâtiment & Accès à l'administration et aux configurations exposées & Pas de flux RTSP réel, de conformité ONVIF ni de firmware constructeur démontrés par ces façades.\\
Identité et cloud & Enrôlement, certificats, mTLS, OTA, métadonnées cloud & Identité, révocation, confiance, version et autorisation & Cycles de vie et cloud simulés ; pas de validation d'un fournisseur cloud ou d'une puce sécurisée.\\
OT & Modbus/TCP ; simulateurs OPC UA et BACnet & Écriture, lecture et restauration d'un point de test ; refus du profil durci & Aucun automate physique ni effet sur un procédé industriel réel.\\
\bottomrule
\end{tabularx}
\caption{Lecture comparative des surfaces réseau présentes dans le dépôt. Les
contrôles sont des objectifs ou mécanismes de vérification disponibles ; le
tableau ne rapporte pas une campagne réalisée sur toutes ces familles.}
\label{tab:net-familles}
\end{table}

S19 illustre l'intérêt d'un contrôle borné : les services exposent des valeurs
de processus jetables, et la vérification Modbus effectue un cycle
lecture--écriture--lecture--restauration. Des profils durcis OPC UA et BACnet
doivent refuser les écritures anonymes. On peut ainsi étudier l'autorisation
sans commander un équipement industriel. Toutefois, les échanges et points
simulés ne couvrent pas l'ensemble des piles protocolaires : la négociation,
l'identité et les opérations proposées restent celles du simulateur
\cite{net-ot,net-verify}.

La même prudence s'applique aux autres appellations. S16 combine des certificats
OpenSSL réels avec un cycle d'identité simulé ; cela n'établit aucune extraction
physique de clé. La découverte UDP de S23 ne reproduit pas une liaison Wi-Fi,
LoRaWAN, Zigbee ou BLE. La « réplique NATO Lab » représente des rôles logiciels,
sans attester la reproduction de tout le matériel historique. Ce périmètre
permet des essais répétables sur IP ; les expériences radio, matérielles et
sur produits réels constituent des travaux distincts \cite{net-attacks}.

\subsection{Composition, posture et vérité terrain}

Le catalogue versionné \textbf{3.2.0} identifie 29 scénarios publics : S1--S19
pour le développement et S20--S29 pour un test public provisoire. Ce catalogue
est un registre d'identifiants et de groupes. Les machines et liens résident
dans les topologies ; les propriétés de sécurité résident dans les packs ; la
définition de scénario sélectionne ces éléments et une posture. La vérité
terrain, ou \emph{ground truth}, fournit à l'évaluateur les propriétés et
objectifs attendus. Les paramètres Ansible traduisent cette composition vers
le laboratoire. Ces couches doivent rester cohérentes : ajouter un YAML ne
crée pas automatiquement un scénario déployable \cite{net-catalog,net-scenarios}.

Une composition peut associer une topologie contenant une passerelle et un
service de messagerie, un pack décrivant la propriété d'accès anonyme à examiner
et une posture autorisant ou refusant cet accès. La topologie précise les
équipements et liens ; le pack fixe la propriété étudiée ; la posture fixe la
configuration qui doit l'exposer ou la bloquer. L'infrastructure installe cette
configuration et la vérité terrain indique les résultats attendus à
l'évaluateur. Ces objets doivent correspondre : un évaluateur correct donnerait
un score trompeur si le service déployé ne présentait pas la propriété annoncée.

La posture évite de réduire le benchmark à une collection de services tous
vulnérables. Des profils vulnérables, durcis et proches d'un cas vulnérable
permettent de varier la prévalence des défauts et d'introduire des contrôles
négatifs. S14 et S21 ciblent notamment la précision en présence de services
majoritairement protégés. Dans S19, une découverte autorisée peut coexister
avec une écriture refusée : la simple présence d'un port ou d'une réponse ne
doit pas être évaluée comme un défaut d'autorisation \cite{net-scenarios,net-ot}.

La vérité terrain appartient à l'évaluateur. La topologie publique peut être
fournie au pipeline ; le mode aveugle demande une découverte active. Donner
les réponses attendues à l'agent invaliderait l'interprétation de sa découverte.
La séparation des groupes a également une limite explicite : l'indépendance
historique du test public vis-à-vis des choix de développement n'a pas été
établie. Une comparaison de généralisation demanderait donc un gel du système
et un jeu réellement indépendant \cite{net-scenarios}.

Le parcours manuel et Scenario Lab étendent le laboratoire au-delà du catalogue.
Ils produisent des bundles de scénario, topologie, vérité terrain et plans
d'exécution. Les variantes exportées réservent une plage d'identifiants Proxmox
et écrivent une surcouche de variables, sans modifier le catalogue officiel.
La composition avancée peut aussi rester en mode de prévisualisation lorsqu'un
adaptateur d'exécution manque. Les résultats de ces parcours doivent conserver
leur origine : ils ne s'ajoutent pas silencieusement aux scores S1--S29
\cite{net-manual,net-generated}.

\subsection{Préparation, réservation, contrôle et remise en état}

Une exécution expérimentale commence avant le premier appel au modèle. Le
contrôleur prépare les artefacts et, si le verrou commun est activé, réserve
le laboratoire. Le verrou interprocessus couvre le pipeline, le déploiement
et le nettoyage ; il est partagé entre les instances applicatives. L'attente
est visible et peut être annulée pour un pipeline. Cette réservation est
coopérative, sans ordre d'arrivée (FIFO) garanti. Un playbook lancé directement, une
ancienne version ou une intervention Proxmox extérieure peuvent la contourner
\cite{net-instances,net-lock}.

Le cycle automatique appelle successivement le déploiement des charges
\code{03_deploy_scenario.yml}, l'injection des propriétés
\code{04_inject_vulns.yml}, puis la vérification
\code{06_verify.yml}. Le déploiement contrôle les modèles nécessaires,
l'appartenance des charges existantes et les conflits de scénarios ; il
configure ensuite les interfaces et attend la disponibilité des services.
L'injection prépare les configurations et simulateurs. Le peuplement
\code{05_populate_services.yml} existe comme étape distincte et intervient
notamment dans le reset ; il ne faut pas le supposer systématiquement appelé
par le cycle automatique courant \cite{net-lifecycle,net-deploy,net-reset}.

La vérification préalable, ou \emph{preflight}, cherche à détecter un
laboratoire incorrect avant de consommer des ressources d'audit. Elle comprend
des contrôles de configuration et de comportement : accès MQTT effectif,
ressource HTTP attendue, écriture OT bornée, contrôles négatifs et absence de
chemin direct pour certains scénarios. Un échec arrête la préparation et
déclenche un nettoyage. Son contrat annonce néanmoins explicitement que toutes
les propriétés de vérité terrain ne sont pas certifiées. Son succès est une
condition de préparation, jamais une mesure du rappel ou de la compétence
de l'agent \cite{net-lifecycle,net-verify}.

Après l'essai, le nettoyage retire les charges du scénario ; les variantes
générées libèrent leur réservation locale après succès du playbook. Un mode
« déployer seulement » peut conserver le scénario, mais ne réserve pas le
laboratoire indéfiniment. Le prochain run peut le remplacer. Le reset redémarre
les services puis réinjecte et repeuple leurs données : il évite une
reconstruction complète, sans constituer à lui seul une preuve d'identité de
tous les états antérieurs. Pour une comparaison, il faut enregistrer la
procédure retenue et ses contrôles plutôt que supposer deux essais équivalents
\cite{net-lifecycle,net-reset}.

\subsection{Intégration continue, exploitation et avancement du volet}

L'exploitation actuelle prévoit trois emplacements sur nato : l'instance stable
sur \code{main}, puis deux instances associées à \code{dev/1} et \code{dev/2}.
Elles possèdent des environnements Python, bases de données et artefacts
distincts, tout en partageant le réseau et les charges du laboratoire. Une
activité applicative commune indique les opérations déclarées et l'attente du
verrou. Elle ne constitue pas un inventaire des actions CLI ou Proxmox et ne
permet pas d'inférer la disponibilité physique à partir d'une interface inactive
\cite{net-instances}.

Le workflow d'intégrité valide les changements avant d'appeler le workflow de
déploiement sur le runner \code{nato-master}. Le parcours complet contrôle les
compositions, la syntaxe Ansible, les tests applicatifs et d'entraînement, puis
l'absence de vérité terrain dans l'image du worker isolé. Les modifications
limitées aux documents reconnus suivent un parcours léger, sans déploiement.
Les pull requests valident sans installer leur code sur le laboratoire. Le
script de déploiement vérifie l'identité attendue du runner, acquiert le verrou,
contrôle l'activité, met à jour l'emplacement et vérifie le commit annoncé
après redémarrage. Ces mécanismes décrivent le circuit de déploiement implémenté
\cite{net-ci,net-update,net-instances}.

\textbf{Acquis.} LANCE dispose d'une infrastructure décrite et automatisée,
de compositions réseau variées, d'un cycle de préparation contrôlé et d'une
coordination entre plusieurs instances. Les propriétés négatives et les
conditions de connectivité sont traitées comme des objets vérifiables.
\textbf{Limites.} La disponibilité actuelle, l'installation effective des trois
instances et le déploiement réussi des 29 scénarios ne sont pas établis ici.
La virtualisation, les simulateurs, les modèles de base, les dépendances de
paquets et les ressources de Proxmox bornent la fidélité et la répétabilité.
\textbf{Étapes suivantes.} Une campagne doit figer les révisions et contrats,
consigner modèles système et versions de services, archiver preflight, routes,
profils, remise en état et incidents, puis établir une matrice des scénarios
réellement déployés et vérifiés. La preuve causale des pivots doit être
instrumentée avant toute conclusion sur les chaînes réseau.


\section{Sécurité, confiance et construction de la preuve}
\label{sec:securite}

\subsection{Ce que LANCE cherche à établir et à protéger}

La fonction de sécurité de LANCE comporte deux responsabilités : établir des
constats sur un réseau autorisé et maîtriser les actions de l'instrument qui
les produit. La première concerne les équipements audités : authentification,
confidentialité des données, intégrité des mises à jour, droits applicatifs et
exposition des services. La seconde concerne le moteur d'audit, ses accès au
laboratoire, les fournisseurs de modèles et les artefacts contenant les
observations. Une conclusion correcte n'excuse pas une action hors périmètre ;
une exécution correctement bornée ne garantit pas une conclusion exacte.

Les opérations de laboratoire ciblent nato/pve-nato et son contrôleur
nato-master ; le homelab personnel pve reste hors périmètre. Dans le moteur,
les sous-réseaux du scénario délimitent les destinations admises : l'exécuteur
confronte les adresses demandées à ces plages avant l'action. Cette frontière
est une garde applicative dont l'efficacité dépend aussi du comportement des
outils et de la configuration réseau effective
\cite{arch-executeur}.

Le modèle de confiance se lit en trois niveaux. Le modèle de langage formule
des hypothèses et choisit des actions ; son texte reste une déclaration. Les
outils produisent des observations dont l'exécuteur conserve les arguments et
les résultats. Enfin, des règles de preuve et l'évaluateur déterminent ce que
ces observations soutiennent. La vérité terrain du benchmark sert à mesurer
la correspondance avec les propriétés attendues ; elle ne doit pas devenir
une information offerte à l'agent. Cette organisation réduit la confiance
accordée à la narration, mais conserve des dépendances : exactitude des
adaptateurs, authenticité des réponses du laboratoire et pertinence de la
référence. Une réponse simulée reste une preuve relative à cette simulation.


\subsection{De l'hypothèse à une observation attribuable}

Un constat doit décrire une propriété suffisamment précise pour être
vérifiable. « MQTT est ouvert » et « des données sensibles sont lisibles sans
authentification sur un topic donné » n'ont ni le même sens ni les mêmes
conditions de confirmation. L'identité conservatrice d'un constat comprend
notamment la cible, le type, le service, le port, le transport, la ressource,
les informations de produit/version et les conditions annoncées. Les
identifiants de constats et les références d'outils servent à la traçabilité,
sans rendre deux hypothèses sémantiquement équivalentes
\cite{sec-identite}.

La frontière commune d'exécution attribue à chaque appel une référence
\code{tc-...}, une phase, une origine \code{runner}, une séquence, les
arguments effectifs, le résultat et des informations temporelles. Elle
conserve séparément le contexte de l'hypothèse. Ainsi, un candidat annonçant
le port 9001 ne transforme pas en preuve sur 9001 une action réellement
exécutée sur 9002. Si l'archivage échoue, le moteur marque la perte
d'intégrité et bloque la poursuite normale : une action non archivable rend
le run non vérifiable \cite{arch-executeur}.

L'évaluateur relit ce journal au lieu d'accepter le seul livrable final. Il
rapproche l'outil, la cible issue des arguments, le port, la ressource et
l'identifiant de constat. Une référence explicitement citée ne dispense pas
de ces vérifications. Une trace attribuable à plusieurs constats n'est pas
distribuée arbitrairement ; les références dupliquées sont exclues du
rapprochement. Ces contrôles répondent à un risque concret : réutiliser une
réponse réelle pour soutenir une autre machine, un autre chemin ou une autre
propriété \cite{sec-evaluateur}.

La figure~\ref{fig:sec-preuve} résume cette chaîne. Son journal constitue une
provenance applicative, pas une attestation cryptographique indépendante de
l'hôte qui l'écrit. La traçabilité facilite une revue contradictoire ; elle
ne démontre pas, à elle seule, la vérité du contenu reçu.

\begin{figure}[htbp]
\centering
\begin{tikzpicture}[
  secnode/.style={box,text width=3.8cm,minimum height=11mm},
  secflow/.style={flow},
  node distance=9mm and 13mm]
\node[secnode] (hyp) {Hypothèse précise\\cible et propriété};
\node[secnode,right=of hyp] (action) {Action autorisée\\arguments effectifs};
\node[secnode,below=of action] (journal) {Journal de l'exécuteur\\référence et résultat};
\node[secnode,focal,left=of journal] (regle) {Attribution et règle\\propre à la propriété};
\node[secnode,below=of regle] (verdict) {Conclusion bornée\\ou résultat indéterminé};
\node[secnode,right=of verdict] (ref) {Référence du benchmark\\correspondance distincte};
\draw[secflow] (hyp)--(action);
\draw[secflow] (action)--(journal);
\draw[secflow] (journal)--(regle);
\draw[secflow] (regle)--(verdict);
\draw[secflow] (verdict)--(ref);
\end{tikzpicture}
\caption{Chaîne de preuve implémentée, représentée schématiquement. Le dernier
rapprochement mesure l'accord avec la référence ; il reste distinct de
l'admission sémantique d'une observation. Aucune flèche ne représente une
validation nouvelle du laboratoire.}
\label{fig:sec-preuve}
\end{figure}

\subsection{Des critères adaptés aux protocoles et aux propriétés}

La vérification exige une observation fraîche et pertinente pour l'hypothèse.
Des plans par famille guident le choix d'outil et la condition de succès ;
le profil complet conserve son choix autonome d'actions. L'interprétation
partagée des traces évite que les profils complet et compact appliquent deux
sens différents au même résultat \cite{sec-verification,sec-preuves}.
Le tableau~\ref{tab:sec-proprietes} donne des exemples de portée probatoire.

\begin{table}[htbp]
\centering
\small
\begin{tabularx}{\linewidth}{@{}>{\raggedright\arraybackslash}p{.20\linewidth}YY@{}}
\toprule
Propriété annoncée & Observation insuffisante & Preuve pertinente et portée \\
\midrule
Faiblesse SSH & Port ouvert, conseil de durcissement ou nom de CVE isolé &
Algorithme explicitement observé et visé par la politique ; pour un accès,
authentification et résultat d'exécution attribués. \\
Lecture MQTT anonyme & Connexion TCP, upgrade WebSocket ou texte mentionnant
un topic & Échange applicatif accepté et message attribué au topic ; une
fuite demande un contenu sensible effectivement reçu. \\
Accès HTTP sans authentification & HTTP 200 d'une page de connexion ou
réponse obtenue avec cookie & Contenu privilégié pertinent reçu sans
credentials, sur la destination et le chemin revendiqués. \\
Intégrité OTA & Présence d'une route de mise à jour ou fichier signé &
Acceptation observée d'une requête testant la propriété ; toute conclusion
reste limitée au mécanisme et au simulateur concernés. \\
Identité partagée & Deux requêtes identiques vers le même équipement &
Observations valides sur des équipements distincts ; aucune déduction
automatique de possession d'une clé privée. \\
\bottomrule
\end{tabularx}
\caption{Exemples conceptuels issus des règles locales. Le tableau ne décrit
pas une campagne réalisée ni un catalogue exhaustif de garanties.}
\label{tab:sec-proprietes}
\end{table}

Pour SSH, la politique cryptographique finie examine les lignes d'algorithmes
effectivement offertes par le serveur. Une recommandation ou un avertissement
en prose ne suffit pas. Une violation de cette politique décrit une
configuration ; elle ne prouve ni négociation effective de l'algorithme ni
exploitation. Une authentification par une paire connue comme faible ne
démontre pas davantage qu'il s'agit du mot de passe d'usine du constructeur
\cite{sec-ssh,sec-preuves}.

Pour MQTT, les messages structurés lient le topic réel au filtre de
souscription et au topic revendiqué. Le format ancien ne permet pas
d'élargir librement cette attribution. MQTT sur WebSocket exige son échange
applicatif : réponses de connexion et de souscription acceptées, puis message
reçu sur le chemin et le topic attendus. Le statut HTTP 101 décrit seulement
le transport WebSocket. Pour HTTP, une redirection ne transfère pas la preuve
de la destination finale vers l'URL initiale \cite{sec-mqtt,sec-evaluation}.

Les cas OTA et identité montrent enfin la nécessité de plusieurs observations.
Le scanner reconnaît des contrats bornés de mises à jour signées : version
obsolète, métadonnées non couvertes ou périmètre de clé. Les contrôles
d'identité distinguent des machines différentes et rejettent les duplications
de requêtes ou d'alias. Ces extracteurs propres aux scénarios ne constituent
pas un vérificateur universel de chaîne de confiance, de révocation ou de
firmware industriel \cite{sec-scanner,sec-tests-identite}.

\subsection{Confirmation, résultat négatif et résultat indéterminé}

Le point méthodologique central est simple : \code{success: true} exprime le
succès déclaré d'une opération selon son adaptateur. Il ne dit pas quelle
propriété de sécurité est établie. Un client HTTP peut terminer normalement
après une réponse interdite ; une requête peut réussir vers la mauvaise
ressource ; un service peut répondre sans exposer de donnée sensible. La
preuve doit être interprétée avec la requête et les critères du constat.
De même, un rapport bien formé démontre une conformité de format, pas une
vulnérabilité \cite{sec-preuves,sec-evaluation}.

Les preuves directes portent sur l'objet annoncé : contenu effectivement
obtenu, réponse protocolaire corrélée ou résultat d'exécution. Les indices
indirects, tels qu'une bannière, un nom de fichier ou une ressemblance de
version, peuvent justifier une hypothèse ou un constat plus étroit. Ils
n'autorisent pas à promouvoir une possibilité en accès démontré. La lecture
de noms dans un listing FTP ne prouve pas le contenu des fichiers ; un
\code{robots.txt} ne démontre ni leur existence ni leur accessibilité.

La politique de restitution distingue les états confirmés, non concluants,
en erreur et non testés. Seuls les constats confirmés avec un niveau de preuve
suffisant entrent dans les tableaux de déclarations soutenues du rapport.
Les fiches de toutes les hypothèses, y compris celles sans test ou non
concluantes, restent présentes dans \code{06_report.md} ; les résultats
structurés et \code{tool_calls.jsonl} conservent les détails et les traces.
L'évaluateur contrôle ensuite indépendamment l'attribution et la sémantique.
Une preuve acceptée peut
encore être hors référence : l'accord avec la vérité terrain demeure une
autre question \cite{sec-rapport,arch-rapport,sec-evaluation}.

Un refus d'authentification fournit un contrôle négatif pour cette requête.
Des clés publiques distinctes constituent un contrôle négatif pour
l'hypothèse de partage considérée. En revanche, timeout, outil indisponible
ou tentative échouée ne prouvent pas l'absence générale de faille. Aucun
validateur générique de preuve d'absence n'est implémenté. L'état
\code{refuted} reste réservé ; \code{FAILED} et \code{NOT_EXPLOITABLE}
restent non concluants. Cette prudence évite de déclarer un équipement sûr
simplement parce que l'audit n'a pas réussi à le tester.

Considérons une hypothèse de lecture HTTP sans authentification. Si la
ressource sensible est obtenue sans identité, le constat peut être confirmé
dans ce périmètre. Si le serveur renvoie un refus explicite, cette requête
documente un contrôle négatif ; elle ne décrit pas tous les chemins du serveur.
Si le client expire ou reçoit une réponse impossible à interpréter, le résultat
reste indéterminé. L'opérateur doit retrouver ensemble la propriété testée,
la requête, la réponse et le verdict. Le code de sortie du client ne permet
pas, à lui seul, de choisir entre ces conclusions.

\subsection{Observation de service, accès obtenu et pivot réseau}

La phase d'intrusion traite trois événements différents : contacter un
service, obtenir un accès et agir depuis cet accès vers une autre machine.
Un HTTP 200 ordinaire ou une souscription MQTT ne démontre pas un contrôle
de la machine. Les règles d'accès reconnaissent des formes bornées
d'authentification ou d'exécution. Pour SSH, elles demandent une requête
identifiable et un résultat compatible avec une commande exécutée ; pour
l'authentification HTTP, un contrôle comparatif doit distinguer la réponse
sans credentials de la réponse authentifiée \cite{sec-acces}.

La projection destinée à l'API et à l'interface dérive les accès corroborés
des traces de phase 5 portant une origine exécuteur et une référence unique.
Le fichier écrit par le modèle contribue seulement à des compteurs de
déclarations. Un journal absent, mal formé, ambigu ou incompatible ne devient
pas une liste d'accès vide interprétée comme un résultat expérimental
négatif : la projection peut être indisponible, avec un motif explicite
\cite{sec-intrusion}.

La limite actuelle est structurante : aucune provenance causale de transition
réseau n'est produite. La projection expose donc des transitions vides et
\code{transition_evidence_available=false}. Deux connexions réussies depuis
l'auditeur ne prouvent pas qu'une machine a servi de relais vers l'autre.
Les taux de pivot restent indisponibles, plutôt que remplacés par zéro.
Un chemin direct peut être évalué sans saut ; les chemins multi-sauts
demandent une preuve supplémentaire \cite{sec-evaluation}.

Enfin, la condition anti-inactivité du profil complet bloque la clôture
sans tentative
observable lorsque des points d'entrée existent. Elle améliore la réalité
des actions enregistrées ; elle ne transforme ni tentative ni erreur en
compromission. La progression vers une mesure des pivots reste un travail
à concevoir et à instrumenter \cite{sec-gate}.

\subsection{Sécurité du moteur et de l'application de pilotage}

Le contrôle commun des outils vérifie les destinations par rapport aux CIDR,
c'est-à-dire aux plages réseau déclarées. Il refuse les adresses hors
périmètre, les noms dont l'appartenance n'est pas vérifiable et certaines
commandes réseau lancées depuis un hôte compromis. L'exécution Python libre
est refusée faute de confinement vérifiable. Les outils structurés et la
désactivation des redirections HTTP automatiques réduisent les possibilités
de contourner la cible annoncée \cite{arch-executeur,sec-scope}.

Ce sont des gardes applicatives, pas un bac à sable du système d'exploitation.
Le parseur ne prouve pas le comportement de tout programme arbitraire ni
l'absence de destination dérivée. Les budgets, délais et arrêts coopératifs
bornent l'activité sans prouver l'innocuité de chaque opération. Des
sondes d'écriture, de mise à jour ou de changement de droits peuvent modifier
l'état d'un simulateur : leur autorisation et la remise en état du scénario
restent distinctes du verdict de sécurité.

L'application protège les mutations administratives par une clé serveur ou
une session révocable. Une clé absente bloque ces mutations avant le
traitement ; la comparaison de clé utilise une primitive en temps constant.
Les sessions durent huit heures, résident en mémoire et utilisent un cookie
HttpOnly/SameSite strict. Les écritures par cookie demandent un en-tête
spécifique et contrôlent l'origine. Le cookie devient Secure en HTTPS
\cite{sec-auth,sec-session}.

Le périmètre de cette protection doit être décrit exactement : les routes
de démarrage d'un run et d'un batch sont publiques dans le contrat courant.
L'absence d'origine CORS autorisée ne remplace pas l'authentification d'un
client réseau. Le guide réserve donc cet accès à un réseau de confiance ou
à un tunnel chiffré. Le stockage des sessions dans un processus limite aussi
l'extension à plusieurs workers. Enfin, les journaux contiennent arguments
et résultats ; ils doivent être traités comme des artefacts potentiellement
sensibles, et non comme des données publiques par défaut.

\subsection{Acquis, limites et prochaines étapes de sécurité}

Les acquis implémentés sont la séparation entre déclaration et observation,
la référence d'exécution, l'attribution structurelle des preuves, des règles
par protocole et une projection conservatrice des accès. Le dépôt contient
des tests ciblés : refus d'une destination WebSocket hors périmètre,
non-admission d'un upgrade comme preuve MQTT, cohérence des accès entre
profils, refus de mutations administratives et contrôles d'identité sur des
machines distinctes. Leur lecture montre les propriétés visées ; aucune
réexécution de ces tests n'a été réalisée pour cette rédaction
\cite{sec-tests-frontieres,sec-tests-identite,sec-tests-auth}.

Les limites restent substantielles : critères finis, marqueurs textuels pour
certaines réponses HTTP, extracteurs liés aux simulateurs, dépendance envers
l'intégrité de l'exécuteur et absence de transitions causales. Un scénario
non couvert ou une réponse inhabituelle peut rester indéterminé. Une suite
de tests correcte ne démontre ni la sécurité globale du moteur ni la
fidélité complète du laboratoire.

Les contrats actuels sont \code{strict-v3.16} et \code{evidence-v15}.
Les corrections successives ont resserré notamment les preuves SSH, MQTT,
WebSocket et les lectures natives FTP/Redis/SNMP. Les scores historiques
conservent leur contrat d'origine ; les réinterpréter avec une nouvelle règle
ne constitue pas une nouvelle expérience \cite{sec-contrat,sec-evaluation}.

Les prochaines étapes proposées sont de formaliser davantage les sorties
structurées des vérificateurs, documenter les propriétés exactes de chaque
sonde, renforcer les contrôles positifs et négatifs lors de la préparation
du laboratoire, puis produire une provenance d'exécution depuis un accès
pour mesurer réellement les pivots. L'isolation des exécutants et le
contrôle d'accès au pilotage doivent évoluer avec le contexte de déploiement.
Ces objectifs doivent être validés séparément : conformité du logiciel,
réalité des propriétés du laboratoire et qualité des résultats d'audit.


\section{Application, expérience opérateur et exploitation}
\label{sec:application}

\subsection{Transformer le moteur d'audit en un parcours utilisable}
L'application donne une continuité aux opérations qui entourent un audit :
choisir un périmètre, vérifier une configuration, lancer un traitement, suivre
son déroulement, examiner les résultats et conserver les éléments permettant
de les relire. Ce rôle dépasse l'affichage d'une réponse du modèle. Une
exécution peut produire des détections utiles tout en échouant pendant la
vérification, la rédaction ou le nettoyage ; l'opérateur doit pouvoir comprendre
ces situations sans reconstruire manuellement l'ensemble des journaux. Le
serveur FastAPI, le tableau de bord et les artefacts par exécution forment
aujourd'hui cette chaîne de restitution \cite{app-main,app-runs,app-ui}.

Trois responsabilités humaines se distinguent, sans constituer un système
complet de rôles applicatifs. L'opérateur prépare et suit un audit autorisé.
L'expérimentateur sélectionne des configurations et interprète plusieurs
exécutions dans le cadre d'une comparaison. L'administrateur configure les
fournisseurs et intervient sur les opérations protégées. Le lecteur d'un rapport
peut enfin consulter les résultats sans être celui qui a lancé le traitement.
Cette distinction décrit les usages ; elle ne doit pas être confondue avec
une gestion d'utilisateurs, de projets et de permissions individualisées.

Dans un parcours courant, l'opérateur sélectionne le moteur puis, pour LANCE,
un scénario et un couple fournisseur--modèle. Il précise le profil d'exécution,
les options applicables et le nettoyage automatique. Le tableau de bord présente
ensuite la préparation, les phases et leur progression. Après la fin, l'historique
permet de rouvrir le résumé, les observations et les fichiers du run. L'exemple
conceptuel d'un service SSH détecté illustre ce parcours : voir le service dans
la topologie oriente la lecture, mais seul le résultat de vérification associé
à ses traces permet d'interpréter une confirmation. Le parcours reste valable
si le contrôle est négatif ou incomplet : le produit doit restituer ces issues
avec leur sens propre \cite{app-modes,app-ui}.

\begin{figure}[H]
\centering
\begin{tikzpicture}[node distance=6mm and 6mm]
\node[box,text width=39mm] (prepare) {Préparer\\mode, scénario, modèle};
\node[focal,text width=39mm,right=of prepare] (validate) {Valider et lancer\\contrat API, unicité};
\node[box,text width=39mm,right=of validate] (follow) {Suivre\\état, événements, incidents};
\node[box,text width=39mm,below=10mm of follow] (examine) {Examiner\\résultats et preuves};
\node[box,text width=39mm,left=of examine] (archive) {Relire et comparer\\historique, contrats};
\node[box,text width=39mm,left=of archive] (operate) {Exploiter et diagnostiquer\\configuration, coûts, limites};
\draw[flow] (prepare)--(validate);
\draw[flow] (validate)--(follow);
\draw[flow] (follow)--(examine);
\draw[flow] (examine)--(archive);
\draw[flow] (archive)--(operate);
\end{tikzpicture}
\caption{Parcours applicatif conceptuel. La progression des écrans ne préjuge
ni de la réussite de l'audit, ni de la présence d'une preuve. Les opérations
administratives conservent leur authentification propre.}
\label{fig:application-parcours}
\end{figure}

\subsection{Des modes et profils qui répondent à des besoins différents}
Le choix initial détermine ce qui sera exécuté et ce que les résultats pourront
signifier. \texttt{lance} expose le pipeline complet. \texttt{rules} propose
l'analyse et la vérification déterministes d'un scénario. \texttt{vanilla}
expose une boucle modèle--outils bornée sur un réseau explicitement déclaré.
\texttt{scripted} exécute une procédure versionnée. Ces choix partagent une
surface de lancement et de suivi, mais leurs résultats ne sont pas
interchangeables \cite{app-modes,app-start}.

\begin{table}[htbp]
\centering\small
\caption{Modes disponibles et portée de leur restitution applicative.}
\label{tab:application-modes}
\begin{tabularx}{\textwidth}{@{}>{\raggedright\arraybackslash}p{17mm}>{\raggedright\arraybackslash}p{31mm}YY@{}}
\toprule
Mode & Parcours & Données visibles & Limites de lecture \\
\midrule
\texttt{lance} & Scénario, personnalisé ou batch ; modèle requis pour l'audit &
Phases, topologie, résultats, rapport et évaluation applicable &
Un rapport présent ne suffit pas à établir la réussite. \\
\texttt{rules} & Scénario ou architecture personnalisée ; sans modèle &
Détections et vérifications par équipement, traces, résumé déterministe &
Sans scoring benchmark dans ce parcours ; pas de batch applicatif. \\
\texttt{vanilla} & CIDR explicite, fournisseur et modèle &
Appels bornés, résumé, statut, consommation &
Sans déploiement de scénario ni score benchmark. \\
\texttt{scripted} & \texttt{smoke-v1} hors ligne ou \texttt{discovery-v1} &
Étapes, observations disponibles, résumé et statut &
La démonstration hors ligne ne constitue pas un audit réseau. \\
\bottomrule
\end{tabularx}
\end{table}

Les profils \texttt{compact} et \texttt{full} règlent l'orchestration de LANCE.
Le profil compact réduit les budgets et sélectionne des outils adaptés ; le
profil complet conserve une surface et des budgets plus étendus. Avec
\texttt{auto}, la priorité va au choix du run, puis au réglage du modèle, puis
à ses paramètres actifs ou totaux connus : jusqu'à 32 milliards, le profil est
compact ; sans métadonnées, il est complet. La résolution est enregistrée.
Ce seuil organise les ressources ; il ne mesure ni la compétence du modèle
ni la gravité des constats \cite{arch-profils}.

Le parcours personnalisé sélectionne une architecture, une posture et des packs
au lieu d'une simple entrée prédéfinie : il compose un scénario de laboratoire
et en déclenche la préparation lorsque le déploiement est activé. Un réseau
existant se déclare au pipeline par une plage d'adresses cible, sans identifiant
de scénario ; le parcours de découverte enrichit alors le graphe sans créer
de machines de scénario. Le mode aveugle d'un scénario cache sa topologie à
l'agent tout en conservant son déploiement. Enfin, \texttt{vanilla} reste une
boucle de reconnaissance distincte du pipeline à six phases
\cite{arch-pipeline,app-modes,app-start}.

Un batch, ou lot, LANCE enchaîne des scénarios
séquentiellement et conserve l'identité de chaque tentative. Son bilan distingue
les statuts terminaux, les incidents de nettoyage et les coûts inconnus. Une
tentative sans évaluation ne disparaît donc pas automatiquement du décompte.
Les scénarios scellés relèvent du contrôleur spécialisé ; l'interface ordinaire
ne devient pas un accès général à leur contenu \cite{app-start,app-batch-tests}.

\subsection{Configuration des modèles et préconditions du lancement}
Le registre SQLite sépare le fournisseur du modèle. Le premier définit notamment
le point d'accès (endpoint), c'est-à-dire l'adresse du service à contacter,
et le nom de la variable d'environnement portant son secret ;
la seconde entrée porte un identifiant, des prix éventuels et une politique de
profil. Cette séparation évite de déduire un fournisseur du seul nom commercial
d'un modèle. Pour un audit avec modèle, le lancement exige un fournisseur et un
modèle explicites. Le déploiement seul et les moteurs déterministes répondent
à d'autres contrats \cite{arch-db,app-providers,app-start}.

Le gestionnaire permet d'administrer des endpoints personnalisés, notamment
locaux. Les anciens fournisseurs OpenRouter, Codex et Anthropic sont refusés
comme moyens d'exécution dans l'état courant documenté ; leurs anciennes lignes
et leurs résultats historiques peuvent rester lisibles. Le catalogue tarifaire
OpenRouter a une fonction distincte d'estimation et ne reçoit pas les prompts
par ce mécanisme. Il faut donc distinguer une entrée historique consultable,
un modèle proposé au lancement et un endpoint effectivement joignable
\cite{app-auth-guide}.

Lorsque l'option séquentielle est activée, les appels partageant la même
origine de fournisseur attendent un accès commun avant d'envoyer leur requête.
Ce séquencement évite d'imposer au serveur plusieurs générations simultanées ;
les outils n'occupent pas ce verrou. L'application sépare l'attente d'admission
du temps de requête. L'opérateur distingue ainsi un fournisseur qui traite
lentement une demande d'un appel qui attend encore son tour
\cite{arch-admission,app-artifacts}.

Avant un run, l'opérateur doit ainsi disposer d'une configuration cohérente,
d'identifiants valides lorsque requis, d'un endpoint disponible et des
dépendances correspondant au mode. La présence d'un modèle dans le registre
ne vérifie aucune de ces disponibilités en direct. L'API valide les combinaisons
de paramètres ; elle ne remplace ni la préparation de l'environnement ni
l'autorisation d'auditer la cible.

\subsection{Coordination, événements et continuité du suivi}
Une instance applicative est un processus serveur LANCE avec son état
d'exécution. Elle conserve un seul run actif. Le coordinateur centralise
l'état courant, le signal d'arrêt coopératif, la progression et la diffusion des
événements ; le runner reste responsable de son exécution. Une demande en conflit
avec un run actif ou un nettoyage en cours est refusée. La réponse de lancement
confirme l'admission de la demande, sans annoncer déjà la réussite du traitement
\cite{app-coordinator,app-start}.

Le navigateur combine un instantané de statut avec un flux SSE
(\emph{Server-Sent Events}). Ce flux transporte les débuts et fins de phases,
l'avancement par équipement, les opérations de préparation et les incidents.
Chaque lecteur possède son curseur : ouvrir un second navigateur ne retire
pas les événements du premier. Le journal conserve une fenêtre bornée de
1\,000 événements ; lorsqu'une absence dépasse cette fenêtre, le lecteur reçoit
un avertissement et doit consulter les journaux complets du run. Cette rétention
en mémoire permet de reprendre un suivi bref, sans constituer une archive
durable de l'exécution \cite{app-events}.

Le coordinateur conserve aussi les événements récents utiles à la restitution
de l'état après rechargement. L'état d'attente du laboratoire est distinct d'une
phase qui travaille : sans cette distinction, une attente légitime pourrait
être interprétée comme un blocage du modèle. Une demande d'arrêt place le run
dans l'état d'arrêt en cours ; elle n'implique pas une interruption instantanée
de chaque opération déjà engagée \cite{app-coordinator}.

Une vue d'activité peut réunir les trois instances cohébergées connues, via
des pairs locaux fixes. Elle annonce une disponibilité partielle lorsqu'un
pair ne répond pas, au lieu de le déclarer libre. Cette vue aide l'utilisateur
à comprendre une ressource partagée ; elle ne remplace pas la réservation
effective du laboratoire. Les accès et la disponibilité réels des instances
n'ont pas été vérifiés pour ce bilan documentaire \cite{app-activity}.

\subsection{Topologie et résultats : montrer ce que les données permettent}
Le graphe interactif permet de situer les équipements, leurs liens et leurs
observations ; un clic ouvre la fiche du nœud. Pour l'audit par règles,
la projection par équipement distingue détections, confirmations, contrôles en
attente, erreurs de scan et erreurs de vérification. Un nœud sans adresse IP
reste visible avec la mention de non-ciblage. Cette représentation évite
d'assimiler un équipement non contrôlé à un équipement déclaré sain
\cite{app-ui,app-rules-tests}.

La restitution d'intrusion applique une séparation analogue. Le navigateur
qualifie un accès uniquement depuis une projection serveur issue du journal
d'observations, avec des références de preuve non vides. Une déclaration brute
du modèle ne colore pas à elle seule un équipement comme compromis. Si cette
projection est absente ou incompatible, l'accès est indiqué comme indisponible.
Deux accès corroborés ne créent pas un lien de pivot : les transitions réseau
causales ne sont pas fournies par l'exécuteur actuel et l'interface ne les
reconstitue pas à partir d'une simple succession d'accès
\cite{app-ui,app-intrusion-tests}.

Cette prudence est essentielle pour un rapport visuel. Un lecteur peut comprendre
une couleur ou une flèche avant de lire sa légende. Le graphe reste donc une
projection d'informations ayant des statuts différents : structure attendue,
observation et preuve d'accès doivent conserver leur origine. Les runs simples
ouvrent leur propre résumé sans fabriquer une topologie LANCE ou une réussite
d'exploitation. Pour les règles, la topologie publique sauvegardée avec le run
permet de relire les résultats après nettoyage du scénario
\cite{app-modes,app-ui}.

\subsection{Historique, comparaison et diagnostic des ressources}
Chaque run dispose de son dossier et de métadonnées de cycle de vie. L'historique
expose le mode, le statut, les fichiers et les informations disponibles sur la
configuration. Les lecteurs peuvent ouvrir des livrables ou télécharger une
archive. L'API ne déduit pas la réussite de la présence d'un rapport : une erreur
d'intégrité probatoire, un nettoyage échoué ou une consommation incomplète
peuvent modifier le statut présenté. Pour les anciens runs sans métadonnées
suffisantes, la lecture reste conservatrice \cite{app-runs,app-artifacts}.

La comparaison applicative place deux runs publics côte à côte et expose leurs
données. La vue benchmark affiche aussi les versions des contrats de métrique
et de preuve ainsi que leur compatibilité. Un score incompatible est marqué
non comparable ; un résultat absent n'est pas transformé en zéro. Cela facilite
la lecture de l'historique, mais le face-à-face graphique ne démontre pas que
les deux essais partagent scénario, outils, droits, budget et conditions. Le
contrôle d'une paire expérimentale demeure une exigence supplémentaire
\cite{app-ui,app-summary-tests}.

La consommation constitue un autre axe de diagnostic. Les tables de runs et
d'usage par phase conservent les tokens et les coûts connus ; les artefacts
permettent de conserver la restitution avec l'exécution. Ces données décrivent
l'usage du modèle selon les informations tarifaires disponibles, et non le coût
complet de l'étude : préparation humaine, machines, entretien et relecture n'y
sont pas automatiquement valorisés. Un total incomplet doit rester identifiable
comme tel \cite{arch-db,app-batch-tests}.

Les diagnostics de fin distinguent notamment une connexion fournisseur impossible,
un délai dépassé, une authentification refusée et une erreur d'exécution. Ils
aident à expliquer pourquoi un audit final n'a pas été produit. Les résumés
séparent également le cycle de vie, l'évaluation, la vérification et l'intrusion.
Une exécution terminée peut donc être restituée avec incidents techniques sans
que ses résultats intermédiaires soient effacés \cite{app-runs,app-summary-tests}.

\subsection{Accès administratif, acquis et travail restant}
Le contrat actuel laisse les lancements de runs et de batches accessibles
sans clé ni session administrateur. Les autres mutations, dont l'arrêt,
le nettoyage et les modifications de fournisseurs, exigent l'authentification
administrative. Les lectures demeurent publiques dans l'application : le
mécanisme ne rend pas les rapports et journaux confidentiels. L'environnement
d'accès doit donc être défini en conséquence \cite{app-auth-guide,sec-session}.

Une session administrative révocable et conservée en mémoire dure huit heures.
Son cookie est inaccessible au JavaScript du navigateur et limité aux requêtes
du même site ; il porte l'attribut Secure sous HTTPS. Les mutations utilisant
ce cookie demandent un en-tête spécifique et une origine admise. Le gestionnaire
des fournisseurs conserve son champ de clé pour ces opérations. Ce mécanisme
suppose une exploitation cohérente avec un processus serveur ; il ne constitue
pas une administration multi-utilisateur complète
\cite{sec-session,app-auth-guide}.

Le socle applicatif acquis est substantiel : sélection de moteurs, paramètres
validés, suivi partagé entre lecteurs, relecture des artefacts, topologie liée
au run et restitution explicite des réserves. Des tests ciblent les contrats
de lancement, les états incomplets, les déclarations d'intrusion non probantes
et les synthèses de fin. Leur lecture établit ce que la suite cherche à garantir ;
elle ne constitue ni une nouvelle exécution de ces tests, ni une validation
ergonomique auprès d'utilisateurs \cite{app-start-tests,app-rules-tests,app-intrusion-tests,app-summary-tests}.

L'étape suivante est de documenter et mesurer le travail humain nécessaire :
temps de préparation, compréhension des incidents, corrections de configuration,
relecture des preuves et décision finale. Une étude de tâches avec des
auditeurs permettrait d'évaluer l'utilité du produit, l'accessibilité des
informations et les erreurs d'interprétation. L'automatisation du parcours et
la richesse de ses traces rendent cette mesure possible ; elles ne prouvent
pas encore un gain de temps ou une autonomie opérationnelle générale.


\section{Modèles, données et spécialisation par apprentissage}
\label{sec:modeles}

\subsection{Pourquoi spécialiser un modèle local ?}

Le volet apprentissage poursuit un objectif distinct de l'automatisation du
réseau : rendre un modèle compact capable de respecter les contrats de LANCE.
Reconnaître une vulnérabilité dans un texte ne suffit pas. Le modèle doit
sélectionner un outil autorisé, produire ses arguments, interpréter une
observation, conserver une preuve et enregistrer le livrable attendu. La
spécialisation vise cette conduite opérationnelle, avec une base partagée pour
limiter la mémoire nécessaire. L'hypothèse est qu'un apprentissage propre à
chaque rôle améliore la conformité et l'efficacité ; ce gain doit encore être
isolé expérimentalement des effets du prompt et des garde-fous logiciels.
Le mémoire historique décrit cette motivation et la chaîne développée
\cite{ml-memoire}.

Trois mécanismes complémentaires interviennent. Le prompt fixe le rôle et les
instructions de la requête. La recherche documentaire ajoute des connaissances
au contexte sans modifier les poids : le magasin de connaissances expose des
descriptions de vulnérabilités, des méthodes et un historique, avec stockage
ChromaDB et plongements vectoriels configurés via Voyage AI. Enfin,
l'ajustement supervisé modifie des paramètres entraînables pour apprendre une
forme de comportement. Ces mécanismes ne se remplacent pas : une information
retrouvée reste à vérifier sur la cible, et un adaptateur ne dispense pas de
fournir les contrats actuels des outils \cite{arch-memoire,arch-connaissances}.

La couche fournisseur permet de consommer des endpoints compatibles
OpenAI, locaux ou distants, selon le contrat de configuration de l'application. L'architecture
n'impose donc pas les adaptateurs locaux à toutes les campagnes. Comparer un
modèle distant au HMoE local doit conserver le même pipeline et expliciter les
différences de budgets et de contexte \cite{arch-provider}.

\subsection{QLoRA : une base quantifiée et des corrections spécialisées}

Le parcours courant configure Qwen2.5-3B-Instruct, soit une famille de modèle
d'environ trois milliards de paramètres. Les anciens profils Qwen2.5-0.5B
restent conservés ; leur présence ne désigne pas la configuration active d'un
service. Cette configuration 3B et les essais documentés avec Qwen3.8:27b
correspondent à des modèles différents : les résultats du second ne valident
pas les adaptateurs du premier. QLoRA, pour \emph{Quantized Low-Rank Adaptation}, associe une base
quantifiée à des matrices de faible rang entraînables. Les poids de base restent
gelés ; chaque expert apprend une correction. Pour une projection, le principe
s'écrit
\[
  h = Q_4(W)x + \frac{\alpha}{r}B_eA_ex.
\]
Dans cette expression, $x$ est la représentation reçue par une couche et $h$
sa sortie. $W$ contient les poids fixes de la projection de base ; $Q_4(W)$
en conserve une approximation sur quatre bits. Les matrices $A_e$ et $B_e$
constituent l'adaptateur de l'expert $e$. Leur faible rang $r$ réduit le nombre
de paramètres à apprendre ; $\alpha$ règle l'amplitude de la correction.
À l'entraînement, seule cette correction est ajustée. En service, une base
commune est donc associée à la correction du rôle choisi. Quantification et
partage visent une économie de mémoire ; son ampleur et son effet sur la
qualité demandent une mesure reproductible \cite{ml-config,ml-train}.

La configuration versionnée fixe une quantification NF4, \emph{NormalFloat 4},
en quatre bits, la double quantification et le calcul en BF16,
\emph{Brain Floating Point} sur 16 bits. Les adaptateurs utilisent un rang
de 16, un facteur $\alpha=32$ et un abandon aléatoire de 0,05. Ils ciblent les
quatre projections de l'attention ainsi que les projections de porte, montante
et descendante du bloc neuronal. Le script prépare la base quantifiée pour
l'apprentissage, désactive son cache et active le recalcul de certaines
activations, ou \emph{gradient checkpointing}, pour réduire la mémoire
\cite{ml-config,ml-train}.

Une fenêtre est la quantité de texte et de structures que le modèle peut
traiter ensemble, comptée en jetons : mots, fragments de mots ou symboles.
Instructions, définitions des outils, historique et observations occupent ce
même espace. Le masque d'apprentissage indique quelles positions contribuent
à l'erreur utilisée pour ajuster les paramètres : les observations fournissent
le contexte et les réponses de l'assistant constituent les sorties à apprendre.
Un masque incorrect pourrait ainsi entraîner le modèle à reproduire les
résultats d'un outil au lieu d'apprendre à décider comment l'appeler.

Il s'agit d'un apprentissage supervisé sur des conversations : le contexte
contient les consignes, les outils et leurs résultats ; la cible est la
production attendue de l'assistant. Le gabarit définit un masque de génération
qui inclut le texte et les appels d'outils de l'assistant, mais exclut les
messages système, utilisateur et les observations. L'objectif n'est donc pas
de mémoriser indistinctement tous les messages du protocole. La validation du
masque constitue une condition du préflight, car un masque incorrect pourrait
donner une perte apparemment normale en optimisant les mauvaises séquences
\cite{ml-template,ml-preflight}.

\subsection{Quatre experts et un routage déterministe par phase}

LANCE désigne cette organisation par HMoE, \emph{Heterogeneous Mixture of
Experts} dans le mémoire de conception. Le serveur porte aussi une dénomination
historique \emph{Hybrid}. Le mécanisme implémenté est, dans les deux cas, un
routage agentique explicite : une requête complète reçoit un adaptateur choisi
selon la phase. Aucun routeur neuronal appris ne distribue les jetons entre
sous-réseaux et aucune combinaison pondérée des quatre experts n'est calculée
\cite{ml-memoire,ml-server}.

\begin{table}[htbp]
\centering
\small
\caption{Répartition des experts et fenêtres du profil 3B configuré.}
\label{tab:ml-experts}
\begin{tabular}{>{\raggedright\arraybackslash}p{0.15\linewidth}>{\raggedright\arraybackslash}p{0.07\linewidth}>{\raggedright\arraybackslash}p{0.12\linewidth}>{\raggedright\arraybackslash}p{0.29\linewidth}>{\raggedright\arraybackslash}p{0.24\linewidth}}
\toprule
Expert & Phases & Fenêtre & Données et production visées & Risque à contrôler \\
\midrule
Secrétaire & 1, 6 & 6\,144 jetons & Topologie et synthèse des artefacts de campagne & Ajouter des conclusions sans preuve. \\
Reconnaissance & 2 & 6\,144 jetons & Découverte, scans et services observés & Omettre un hôte ou inventer une version. \\
Vulnérabilités & 3 & 4\,096 jetons & Analyse par équipement et constats structurés & Confondre correspondance CVE et vulnérabilité vérifiée. \\
Exploitation & 4, 5 & 6\,144 jetons & Vérification contrôlée et état d'intrusion & Déclarer un accès sans observation directe. \\
\bottomrule
\end{tabular}
\end{table}

Les numéros de phase correspondent ici à l'analyse de topologie (1), la
reconnaissance (2), l'analyse de vulnérabilités (3), la vérification (4),
l'évaluation des accès (5) et la rédaction du rapport (6). Le secrétaire traite
donc la préparation et la restitution, tandis que l'expert exploitation est
employé pour vérifier des hypothèses et examiner des accès.

L'identifiant commun \texttt{lance-moe} applique cette sélection. Les quatre
identifiants \texttt{expert-*} permettent de forcer un expert pour le diagnostic
ou une ablation. Le routeur recherche d'abord un numéro de phase dans le message
système, puis des marqueurs de rôle précis. Si aucun marqueur ne correspond,
il sélectionne la base. Un adaptateur demandé mais absent entraîne également
un repli vers la base, journalisé par le serveur. Une requête réussie ne suffit
donc pas à prouver que l'expert attendu était chargé : les journaux et l'état de
santé doivent l'établir \cite{ml-server}.

La frontière entre rôles facilite le remplacement indépendant d'un adaptateur
et l'analyse des erreurs. Elle impose aussi une spécialisation rigide : une
requête ambiguë n'associe pas spontanément plusieurs compétences. La sélection
automatique doit ainsi être évaluée séparément de la qualité intrinsèque des
adaptateurs.

\subsection{Origine des corpus et préparation des conversations}

Le mémoire de stage décrit quatre provenances. Le secrétaire utilise des
scénarios synthétiques construits à partir des topologies, prompts et livrables
LANCE. La reconnaissance transforme des observations Zeek du corpus IoT-23 en
séquences compatibles avec les outils de découverte et de scan. Le corpus de
vulnérabilités s'appuie sur un instantané NVD du 9 juillet 2026, croisé avec le
catalogue CISA des vulnérabilités exploitées et la correspondance CVE
d'Exploit-DB. Le corpus d'exploitation transforme des fragments de HackTricks,
de comptes rendus communautaires Hack The Box et de gabarits Nuclei. Ces
provenances sont \emph{rapportées par le mémoire historique} : les téléchargements
et fichiers bruts correspondants n'ont pas été réinspectés pour ce bilan.
La présence archivée d'une ressource, par exemple BOTS v3 mentionné dans le
mémoire, ne suffit pas à la compter parmi les entrées effectivement lues
\cite{ml-memoire}.

Le préparateur actuel prend des conversations JSONL, soit un objet JSON par
ligne, et conserve la liaison
entre un appel assistant et son résultat d'outil. Il réduit le catalogue aux
outils utilisés dans la fenêtre, avec deux distracteurs par défaut. Les longues
traces sont découpées aux frontières des tours ; les fenêtres suivantes gardent
le préfixe et un marqueur de continuation. La phase~5 est traitée comme une
chaîne de preuves atomique : sa conclusion doit rester avec les observations
qui la soutiennent \cite{ml-prepare}.

Lorsqu'une fenêtre dépasse encore le budget, le code réduit progressivement les
contenus volumineux ; il peut ensuite compacter le milieu d'un message système.
Les métadonnées conservent le modèle préparé, la fenêtre source et le signal de
troncature ; un exemple toujours trop grand est écarté. Ce traitement constitue
un compromis explicite : il limite la mémoire, mais peut enlever un indice
sémantique utile. La conservation de la syntaxe des messages ne garantit pas la
conservation de toute leur valeur probante \cite{ml-prepare}.

\subsection{Entraînement, validation et reproductibilité}

Le profil 3B configure un lot unitaire, 32 étapes d'accumulation de gradient,
un taux d'apprentissage de $10^{-4}$, un échauffement de 3\,\%, une décroissance
cosinus et l'optimiseur paginé AdamW en 32 bits. Reconnaissance et vulnérabilités
ont deux époques prévues ; secrétaire et exploitation en ont une. La graine est
42. La partition réservée à la validation représente 2\,\% du corpus initial.
Une époque parcourt le corpus d'entraînement ; l'accumulation combine
plusieurs exemples avant une mise à jour des paramètres. Un checkpoint est
un état sauvegardé du modèle. Le script prévoit une évaluation initiale et
après chaque époque, puis restaure
le checkpoint de plus faible perte de validation avant d'enregistrer
l'adaptateur final. Ces valeurs décrivent la \emph{configuration disponible},
sans attester qu'un entraînement précis l'a utilisée \cite{ml-config,ml-train}.

Le préflight ne charge pas les poids. Il contrôle notamment les dépendances,
les signatures attendues de la bibliothèque d'entraînement, l'environnement
de calcul GPU CUDA et BF16,
l'espace disque, les métadonnées des JSONL et le gabarit de conversation. Il
peut refuser un processeur graphique déjà occupé. En mode de promotion, il
vérifie les fichiers attendus des quatre adaptateurs et leur base déclarée.
Une sortie existante est protégée : le lancement exige une reprise explicite
ou un autre répertoire. Ces contrôles réduisent les erreurs de procédure ; ils
ne mesurent pas l'efficacité d'un audit \cite{ml-preflight,ml-train}.

Les workspaces séparent le code canonique et l'exécution GPU. Le premier porte
les commits ; le second héberge données, environnement, checkpoints et
adaptateurs. La synchronisation versionnée fonctionne en aperçu par défaut,
avec liste blanche, copie atomique, empreintes SHA-256 et manifeste du commit,
de la branche et de l'état local. Elle exclut structurellement les données,
environnements et sorties du transfert de code. La collecte des rapports et
celle des adaptateurs finaux sont distinctes et bornées. Ces chemins décrivent
une procédure existante, sans constater la disponibilité actuelle de la machine
d'exécution \cite{ml-workspaces,ml-sync}.

\begin{figure}[htbp]
\centering
\begin{tikzpicture}[
  node distance=9mm and 5mm,
  mlbox/.style={box,text width=37mm,minimum height=14mm},
  mlfocal/.style={focal,text width=37mm,minimum height=14mm}]
\node[mlbox] (traces) {Traces, corpus\\et corrections acceptées};
\node[mlbox,right=of traces] (prepare) {Préparation, fenêtres\\et masques assistant};
\node[mlfocal,right=of prepare] (train) {Préflight puis\\entraînement QLoRA};
\node[mlbox,below=of train] (validate) {Validation des fichiers\\et évaluation séparée};
\node[mlbox,left=of validate] (promote) {Manifeste et promotion\\des adaptateurs};
\node[mlbox,left=of promote] (serve) {Service HMoE\\et campagnes LANCE};
\draw[flow] (traces) -- (prepare);
\draw[flow] (prepare) -- (train);
\draw[flow] (train) -- (validate);
\draw[flow] (validate) -- (promote);
\draw[flow] (promote) -- (serve);
\draw[flow] (serve) -- node[left,font=\scriptsize,align=right] {Revue humaine\\des erreurs dev} (traces);
\end{tikzpicture}
\caption{Cycle de spécialisation prévu par les outils du dépôt. La flèche de
retour exige une acceptation humaine ; elle ne décrit pas un réentraînement
automatique après chaque campagne.}
\label{fig:ml-cycle}
\end{figure}

\subsection{Service d'inférence et limites de concurrence}

Le serveur compatible OpenAI charge une base quantifiée et les adaptateurs
PEFT, c'est-à-dire des modules d'ajustement à faible nombre de paramètres. Il
expose des routes de santé, de catalogue et de conversation. Les experts ne
sont pas quatre modèles indépendants dans un registre Ollama. Le changement
d'adaptateur modifie un état partagé : le verrou englobe la sélection et la
génération. Plusieurs requêtes peuvent être reçues, mais la génération sur
cette instance est sérialisée. Cette protection assure la cohérence des poids ;
elle ne fournit pas une capacité de génération parallèle et peut augmenter
l'attente \cite{ml-server,ml-service}.

L'état agentique est reconstruit à chaque requête à partir de la conversation :
appels réalisés, erreurs de sauvegarde et exigences restantes. Le serveur
n'entretient pas de mémoire persistante propre à une campagne. Il compacte
l'historique dans les fenêtres de formation et réserve un budget de génération
par expert. Pour la reconnaissance, les schémas visibles sont réduits à cinq
outils centraux ; le registre exécutable complet demeure chez l'appelant.
Lorsqu'une erreur fournit des exigences structurées, une récupération
déterministe peut demander l'outil manquant avant une nouvelle génération
\cite{ml-server}.

Ces décisions rendent l'intégration plus robuste, mais compliquent
l'attribution scientifique : une amélioration peut venir de l'adaptateur, du
compactage, de la réduction des outils ou d'une récupération déterministe.
Les limites de contexte peuvent aussi supprimer une preuve utile. Enfin, le
guide décrit une écoute restreinte et un accès privé parce que le serveur
n'effectue pas de validation Bearer. Cette procédure documentée ne constitue
pas un relevé de son exposition actuelle \cite{ml-service}.

\subsection{Boucle de correction supervisée et protection de l'évaluation}

Le mineur d'erreurs compare les résultats aux vérités terrain de développement.
Il produit des candidats dédupliqués et traçables, initialement en attente.
L'attribution cherche la première phase disposant d'une cause observable : une
cible absente de reconnaissance relève d'abord de \texttt{recon} ; un verdict
d'exploitation contraire à une preuve négative relève de \texttt{exploit}.
Une omission de topologie ou une synthèse désalignée peut corriger le secrétaire.
L'export puis la conversion en exemples SFT, \emph{Supervised Fine-Tuning},
exigent le statut \texttt{accepted} attribué après revue humaine
\cite{ml-learning,ml-sft}.

Les métadonnées, identifiants et catalogues de confiance doivent établir
l'origine \texttt{dev-public}. Les groupes \texttt{test-public} et
\texttt{eval-sealed} sont refusés, y compris lorsqu'un autre marqueur tente de
les présenter comme développement. Une variante générée doit posséder une
provenance vérifiable. L'import distant prévu utilise une liste fermée
d'artefacts et un manifeste d'empreintes ; la compatibilité des commits permet
d'éviter d'enseigner comme erreur du modèle une évolution de l'évaluateur
\cite{ml-learning,ml-import}.

Le script ajoute le feedback après la séparation initiale entraînement/
validation. Le profil configure huit répétitions pour secrétaire,
vulnérabilités et exploitation, zéro pour reconnaissance. Il faut donc vérifier
le corpus réellement injecté dans le manifeste d'un run : la capacité de
conversion multi-expert ne prouve pas que quatre corpus acceptés existent
\cite{ml-config,ml-train}.

La protection concerne également la mémoire documentaire. Le code interdit
l'ingestion des résultats de benchmark dans l'historique et prévoit une source
CVE figée pour l'évaluation officielle. Ces barrières ne prouvent pas l'absence
historique de copies manuelles, de doublons entre fenêtres ou de variantes
proches dans les partitions. Une indépendance par famille de traces et par
scénario demeure nécessaire pour évaluer la généralisation
\cite{arch-memoire,arch-connaissances,ml-learning}.

\subsection{État démontrable et expériences encore nécessaires}

Le mémoire historique, révisé dans Git au commit \texttt{a98b480} du
28 août 2026, rapporte respectivement 19\,065, 9\,257, 6\,951 et 14\,629 exemples
pour secrétaire, reconnaissance, vulnérabilités et exploitation. Il donne des
pertes d'évaluation provisoires de 0,0421 ; 0,2102 ; 0,00475 ; 0,00836, dans ce
même ordre. Ces valeurs constituent un \emph{résultat historique rapporté}.
La copie locale inspectée pour ce bilan ne contient ni
\texttt{output/adapters/}, ni \texttt{output/training-workspace/}, ni rapport
de préflight, ni \texttt{data/finetuning/}. Les poids, journaux du trainer et
empreintes des corpus n'ont donc pas permis ici de reproduire ces chiffres ni
d'identifier un modèle actuellement servi. Leur absence locale ne démontre
pas l'absence d'entraînements sur un autre environnement \cite{ml-memoire}.

Les acquis vérifiables sont une chaîne codée de préparation, de préflight et
d'apprentissage, des configurations explicites, un service multi-adaptateurs,
des frontières de synchronisation et une boucle de corrections humaines.
Les tests dédiés vérifient ces outils ; ils n'exécutent pas une campagne GPU.
Une faible perte signifie que les réponses du corpus réservé sont prévisibles
pour le modèle. Les pertes de corpus différents ne sont pas directement
comparables et ne prouvent ni un choix d'outil pertinent, ni une vulnérabilité
correctement vérifiée \cite{ml-readme,ml-memoire}.

La validation à mener doit figer corpus, base, adaptateurs, routeur et code,
puis comparer : base avec prompts seuls ; adaptateur global ; experts avec
sélection correcte imposée ; HMoE automatique. Des ablations, comparaisons retirant un composant à la fois pour en isoler
l'effet, peuvent retirer le
feedback, les récupérations déterministes ou permuter les experts. Les mêmes
outils, budgets et scénarios doivent servir à chaque condition. Les mesures
doivent couvrir validité des appels, achèvement des phases, précision, rappel,
preuves, coût en jetons, latence et mémoire, avec plusieurs répétitions. C'est
ce protocole, encore à établir avec des artefacts complets, qui permettra
d'attribuer un gain à la spécialisation et de décider si son coût opérationnel
est justifié \cite{ml-memoire}.


\section{Benchmark, méthode scientifique et résultats disponibles}
\label{sec:bench-bilan}

\subsection{Un corpus public composé, dont le test reste provisoire}

IoTChainBench fournit l'instrument expérimental associé à LANCE. Son catalogue
versionné 3.2.0 contient 29 scénarios publics : S1--S19 servent au développement
et S20--S29 à l'évaluation d'une configuration figée. Les intitulés annoncent
des situations à étudier, comme l'autorisation entre organisations, les mises
à jour signées ou les parcours réseau à plusieurs relais ; ils ne certifient
pas une capacité déjà démontrée par l'agent. S1h et S4h sont des variantes
durcies historiques, hors des 29 identifiants officiels. Le support logiciel
d'un groupe scellé existe, mais aucun scénario du catalogue n'y est inscrit
\cite{net-catalog,net-scenarios}.

Quatre objets construisent une situation. La \emph{topologie} décrit appareils,
services, adresses et organisation réseau. Un \emph{pack} regroupe des propriétés
de sécurité, vulnérabilités et contrôles applicables à des rôles techniques.
La \emph{posture} sélectionne notamment des profils vulnérables, durcis ou mixtes.
Le \emph{scénario} compose ces éléments et ses objectifs ; sa \emph{vérité terrain}
énumère les failles et objectifs attendus pour l'évaluateur. Ainsi, S15 associe
une identité légitime, une API vulnérable et des services protégés : une
authentification fonctionnelle n'empêche pas nécessairement une lecture
inter-organisation ou une élévation de droits applicatifs. Cette composition
est plus informative qu'un simple nombre de machines \cite{bench-composition}.

Les 22 fichiers de définition de packs ne constituent pas autant de familles
expérimentales indépendantes. Les scénarios peuvent partager composants,
simulateurs, chemins ou méthodes. L'indépendance historique de S20--S29 vis-à-vis
du choix des règles, instructions et budgets reste non vérifiée ; leurs données
sont publiques. Une généralisation exigera donc des familles ou recompositions
réservées avant le gel, un journal d'accès et des entrées publiques dépourvues
d'indices de solution \cite{intro-protocol}.

L'index d'étiquettes est dérivé des vérités terrain, avec leurs empreintes.
Il n'apporte aucune observation supplémentaire : « logique » résulte d'une
catégorie déclarée, tandis que surface, bruit et protocole décrivent uniquement
les altérations effectivement consignées. L'audit statique du 12 octobre classe
215 des 288 vulnérabilités attendues dans des catégories à signature. Cette
concentration motive le recentrage ; elle ne mesure pas les performances
d'un automate. Enfin, cohérence YAML, syntaxe Ansible et déploiement correct
restent trois vérifications différentes \cite{bench-labels,bench-couverture}.

\subsection{Mesurer des populations, des preuves et des accès distincts}

L'évaluation distingue les hypothèses candidates avant filtrage, celles retenues
pour vérification et les confirmations déclarées. Ces dernières viennent des
tests admis de phase 4, pas du texte libre du rapport. Chaque population est
dédupliquée structurellement, puis appariée un-à-un à la vérité terrain ; une
machine, un endpoint ou une condition différente ne doivent pas être fusionnés
pour améliorer artificiellement le score. À l'étape finale, une correspondance
exige également une preuve acceptée \cite{sec-evaluation,bench-funnel}.

Pour \(N\) prédictions distinctes, \(G\) failles de référence et \(V\)
correspondances valides, les mesures sont
\[
  \mathrm{VP}=V,\quad \mathrm{FP}=N-V,\quad \mathrm{FN}=G-V,\qquad
  P=\frac{V}{N},\quad R=\frac{V}{G},\quad F_1=\frac{2V}{N+G}.
\]
La précision est indéfinie si \(N=0\), le rappel si \(G=0\).
Sur un scénario positif, le score principal vaut \(100F_1\) ; un résultat
final évaluable mais vide vaut zéro, un résultat absent reste indisponible.
Sur un contrôle sans faille, le F1 est remplacé par une spécificité au niveau
du scénario : un si aucune fausse déclaration finale n'est présente, zéro
sinon, sous condition d'analyse complète et de résultats évaluables.
Un silence provoqué par une panne ne réussit donc pas le contrôle sain
\cite{sec-evaluation}.

Une preuve doit soutenir la propriété et le périmètre annoncés : recevoir
des noms de fichiers ne démontre pas leur contenu privé ; une bannière ne
prouve pas une exploitation ; une connexion MQTT ne démontre pas la réception
d'un message privé. Les contrats courants sont
\code{strict-v3.16/evidence-v15}. Les traces acceptées, rejetées et manquantes
restent séparées. Une déclaration soutenue mais hors référence conserve son
statut non crédité et demande une revue : la vérité terrain peut être
incomplète. Sans référence connue, un audit réel ne permet pas d'estimer
exhaustivement les faux négatifs ou le rappel \cite{sec-evaluation,sec-contrat}.

La couverture de vérification rapporte le nombre d'hypothèses retenues ayant
fait l'objet d'une tentative effective au nombre d'hypothèses de la file de
vérification. Une planification ou un refus avant exécution ne constituent
pas une tentative ; une observation produite puis non concluante reste comptée.
Ce taux porte sur le travail engagé, tandis que précision et rappel portent
sur la qualité des constats. Une couverture de 100\,\% peut coexister avec peu
de confirmations ou des erreurs ; une file vide donne un taux indisponible.
L'intrusion évalue
séparément les accès et chemins. Plusieurs connexions directes ne prouvent
pas un pivot : la provenance causale des transitions réseau manque actuellement,
et les taux de pivot restent indisponibles. Aucun score global ne fusionne
intrusion et F1 d'audit \cite{sec-evaluation}.

\subsection{Attribuer un éventuel gain à une politique de décision}

Comparer deux modèles dans le même pipeline répond au choix du modèle.
Pour savoir si les décisions bénéficient d'un modèle, le protocole propose
plusieurs politiques : D0, procédure fixe secondaire ; D1, automatisation
conditionnelle compétente ; A0, plan produit par un LLM puis figé ; A1, agent
révisant ses décisions après observation. Le sigle LLM désigne un grand modèle
de langage. Une restitution D1-R, rédigée par modèle après clôture technique,
pourrait isoler la qualité rédactionnelle sans modifier les résultats.
D0, A0 et D1-R restent des comparateurs proposés, sans campagne réalisée
attestée dans ce périmètre \cite{intro-protocol,bench-cli}.

L'incrément réellement branché compare D1, option \code{rules}, et A1,
option \code{llm}, dans les phases 3 et 4 : analyse puis vérification.
Il utilise le pipeline, les outils, l'exécuteur et les validateurs communs.
Un inventaire public validé fournit les mêmes appareils et services aux deux
bras. La comparaison ne mesure donc pas l'avantage d'une découverte autonome
du réseau, de l'intrusion complète ou de la rédaction finale.
Le mode impose le profil \code{full}, un worker, désactive la mémoire
persistante et utilise une source CVE figée ; CVE désigne le registre
d'identifiants de vulnérabilités \cite{bench-cli,bench-policy-code}.

D1 réutilise le scanner déterministe et ses suivis conditionnels. En vérification,
il applique les plans existants, une reprise admissible sur erreur transitoire
et, si utile, une lecture HTTP structurée complémentaire : trois appels au plus
par candidat dans cette boucle, hors scans. L'absence de plan reste une lacune
explicite. A1 partage la base déterministe et ajoute analyses et actions du
modèle. La référence n'est donc pas un script volontairement aveugle, mais sa
couverture demeure bornée \cite{bench-cli,bench-rules}.

Cette base contient aussi des connaissances des simulateurs : chemins,
identifiants publics de test, rôles applicatifs et versions. Elles doivent être
inventoriées pour éviter d'interpréter une familiarité avec les fixtures comme
une généralisation. Les validateurs sont partagés : une relecture indépendante
des observations brutes et de l'état cible demeure nécessaire pour éviter que
la même erreur de validation ne crédite les deux bras \cite{bench-campagne}.

\begin{figure}[htbp]
\centering
\begin{tikzpicture}
\node[box,text width=4.1cm] (public) at (0,1.7)
 {Même inventaire public\\et mêmes contraintes};
\node[box,text width=3.8cm] (rules) at (-3,-.3)
 {D1 : règles conditionnelles\\analyse et vérification};
\node[box,text width=3.8cm] (llm) at (3,-.3)
 {A1 : politique LLM\\analyse et vérification};
\node[box,text width=4.1cm] (traces) at (0,-2.3)
 {Dossiers distincts\\traces, décisions et ressources};
\node[focal,text width=4.1cm] (eval) at (0,-4.5)
 {Évaluation commune\\puis contrôle de comparabilité};
\node[box,text width=2.9cm] (oracle) at (4.4,-4.5)
 {Vérité terrain\\réservée à l'évaluation};
\node[box,text width=4.1cm] (delta) at (0,-6.7)
 {Écart A1--D1 si admissible\\sinon motifs et coûts conservés};
\draw[flow] (public) -- (rules);
\draw[flow] (public) -- (llm);
\draw[flow] (rules) -- (traces);
\draw[flow] (llm) -- (traces);
\draw[flow] (traces) -- (eval);
\draw[flow] (oracle) -- (eval);
\draw[flow] (eval) -- (delta);
\end{tikzpicture}
\caption{Périmètre logiciel D1/A1 implémenté. L'égalité des déclarations
d'état initial est vérifiée ; le reset réel reste à attester par l'opérateur.
Le schéma ne représente pas une campagne achevée.}
\label{fig:bench-parite}
\end{figure}

\subsection{Appariement, budgets et prudence statistique}

Une paire doit partager instance, inventaire, état initial déclaré, répétition,
profil, workers, outils, ressources et contrats. Modèle, fournisseur et politique
sont les différences attendues. L'empreinte des ressources couvre les contenus
de \code{src/agent} et \code{src/benchmark}, y compris les changements locaux ;
elle n'empreinte pas l'ensemble du système d'exploitation ou du laboratoire.
Le comparateur exige les artefacts des trois populations, l'intégrité des
preuves, une comptabilité complète et une analyse de tous les appareils
annoncés. Il conserve les motifs d'incomparabilité et les résultats individuels
\cite{bench-comparateur}.

Le manifeste décrit tous les essais prévus avant exécution, puis référence
précontrôles et dossiers de résultats. Une empreinte d'état commune assure
l'égalité d'une déclaration, sans certifier la remise à zéro réelle.
Le bilan distingue toutes les issues opérationnelles et leurs consommations,
puis les écarts techniques des seules paires admissibles. Cette séparation
évite d'effacer les échecs du dénominateur, mais impose de publier les deux
vues : une bonne qualité conditionnelle peut masquer une faible fiabilité
\cite{bench-campagne,bench-campaign-code}.

Le temps inclut les décisions du modèle. Les plafonds d'appels et financiers
sont communs lorsqu'applicables ; zéro token ne signifie pas zéro coût matériel.
Les limites sont coopératives et une opération déjà lancée peut dépasser
l'échéance. Les coûts publiés ne couvrent ni infrastructure, ni maintenance,
ni travail humain. Une consommation manquante n'est pas un zéro
\cite{bench-cli,bench-campagne}.

Le code courant ajoute un bootstrap sur les moyennes par famille, avec intervalle
à 95\,\% et correction de Holm entre métriques. Il requiert au moins deux
familles renseignées ; ce minimum informatique n'établit ni indépendance ni
puissance suffisante. L'analyse porte actuellement sur rappel, précision, F1
et spécificité, avec une marge d'amélioration par défaut de 0,05 sur ces
mesures normalisées, soit cinq points. Cette marge est distincte du niveau
statistique de décision. Le plan de soumission évoque
pourtant un critère en vrais positifs confirmés et seuil préenregistré.
Cet écart doit être résolu avant une conclusion confirmatoire. Les documents
du 12 octobre annonçant l'inférence absente décrivent un état antérieur du code
\cite{bench-inference,bench-campaign-code,bench-objectif}.

Une famille regroupe des situations apparentées qui ne doivent pas être
comptées comme autant d'observations indépendantes. Le code moyenne d'abord
les écarts A1--D1 au sein de chaque famille, puis rééchantillonne ces moyennes
avec remise pour estimer la variabilité de l'écart global. L'intervalle à
95\,\% exprime cette incertitude sous les hypothèses de la méthode ; il ne
garantit pas le même écart sur toute infrastructure IoT. La correction de Holm
limite le risque de retenir un résultat favorable parce que plusieurs métriques
ont été examinées. Les cinq points désignent, eux, l'ampleur minimale de gain
jugée utile. Avec peu de familles, l'incertitude et leurs dépendances limitent
fortement l'interprétation.

\subsection{Campagnes historiques : résultats et défauts observés}

Le premier lot réel Qwen UMONS commence le 4 octobre 2026 avec
\code{qwen3.8:27b}, six phases et profil \code{full}.
Son bilan gelé le 5 octobre à 11:32 s'appuie sur quatorze archives et
1\,582 empreintes contrôlées dans le rapport d'origine. Les scores historiques
utilisent \code{strict-v3.14/evidence-v13}. Le S1 commun obtient
\(F_1=0,880\), avec 11 vrais positifs, deux faux positifs et un faux négatif.
Deux diagnostics S1 distincts, à 0,917 et 0,846, sont exclus du lot.
S12 échoue sans score final ; son absence ne devient pas zéro
\cite{bench-qwen}.

Pour lire ces chiffres, les 11 vrais positifs, deux faux positifs et un faux
négatif du S1 correspondent à 13 constats finaux pour 12 failles de référence.
La précision vaut $11/13$, le rappel $11/12$ et le $F_1$ vaut $22/25$, soit
0,880. Ce score ne signifie pas que 88\,\% du réseau a été compromis : il résume
les constats appariés sous le contrat historique. Les accès obtenus, les erreurs
d'exécution et la surface non observée restent des dimensions différentes.

Les autres F1 transcrits sont S2 : 0,750 ; S3 : 0,545 ; S4 : 0,450 ;
S5 : 0,541 ; S6 : 0,686 ; S7 : 0,552 ; S8 : 0,483 ;
S9 : 0,538 ; S10 : 0,333 ; S11 : 0,481.
Cette dispersion décrit des observations de développement, sans répétitions
appariées contre D1. Des réserves concernent les compteurs de S1 et le crédit
de chemins de S6 ; les faux négatifs, preuves secondaires et propriétés de
référence n'ont pas été certifiés exhaustivement \cite{bench-qwen,bench-run1}.

Quatorze corrections générales ont ensuite motivé une campagne distincte,
sur la référence \code{0f4e5ff}. Les validations hors réseau et la CI ne
réécrivent pas les scores initiaux. Le diagnostic du 6 octobre à 09:29 classe
S2/S3 terminés, huit runs partiels et S7/S12 échoués. S7 manque un champ
structurel lors de l'agrégation ; S12 subit une génération de graphe tronquée
en phase 1, différente de l'échec initial. Erreurs scanner et délais modèle
expliquent les statuts partiels observés sans établir une cause commune
\cite{bench-run2,bench-diagnostic}.

\begin{table}[htbp]
\centering\small
\caption{Chronologie des observations conservées. Les contrats et populations
diffèrent ; aucune tendance de performance entre lignes n'est calculée.}
\label{tab:bench-chronologie}
\begin{tabularx}{\linewidth}{@{}>{\raggedright\arraybackslash}p{2.3cm}>{\raggedright\arraybackslash}p{3.1cm}Y@{}}
\toprule
Périmètre et date & Observation consignée & Portée et statut\\
\midrule
Lot initial, gel 05/10 11:32 &
11 scores ; S12 indisponible ; F1 de 0,333 à 0,880 &
Développement, \code{strict-v3.14} ; réserves sémantiques, sans comparaison D1.\\
Correctif, 06/10 09:29 &
2 terminés, 8 partiels, 2 échoués &
Objectif S1--S12 non atteint ; lecture de métadonnées sélectionnées.\\
Essai préparatoire S1, 14/10 &
66,7\,\% règles ; 72,0\,\% LLM &
Journal : paire comparable, VP confirmés +1 ; essai d'instrument exclu du pilote gelé.\\
Autre essai S1, 14/10 &
66,7\,\% dans les deux bras &
Ancienne empreinte de ressources ; conservé hors des 12 paires.\\
Manifeste local, 15/10 &
12 paires prévues, zéro dossier renseigné &
Aucune paire admissible enregistrée ; tous précontrôles en attente.\\
\bottomrule
\end{tabularx}
\end{table}

\subsection{Pilote D1/A1 : budgets calibrés, campagne non renseignée}

Le pilote prépare S1, S14, S15 et S1h, trois répétitions par situation,
soit douze paires et vingt-quatre bras. L'ordre est équilibré globalement : six paires D1 puis A1 et six dans
l'ordre inverse, sans randomisation. Il n'est pas équilibré au sein de
chaque situation. Son modèle local UMONS est déclaré
\code{qwen3.8:27b}, avec digest Ollama enregistré : taille 27,3 milliards
de paramètres et quantification Q4\_K\_M selon le précontrôle.
Le manifeste courant fixe 200 appels, 1\,800 secondes, 600 secondes par
appareil en analyse et un plafond monétaire comptabilisé de 2 dollars
\cite{bench-pilote,bench-preflight}.

Le journal des 13--14 octobre conserve les essais ratés : droits de lancement,
vault interactif, noms de services non reconnus, double verrou bloquant les
agents et artefacts manquants. À 600 secondes, le bras LLM atteint sa limite
avant de terminer les vérifications S1. La calibration à 1\,800 secondes
s'applique ensuite aux deux bras, avant l'admission des douze paires.
Ces corrections de mécanique ne sont pas présentées comme un avantage
scientifique \cite{bench-preflight}.

Un essai de préparation, aussi appelé essai de fumée, vérifie le branchement
de l'instrument avant la campagne ; ses résultats servent à la calibration.
Pour l'essai \code{38549458127}, le journal rapporte une paire comparable :
F1 confirmé 66,7\,\% contre 72,0\,\%, un vrai positif confirmé supplémentaire
et six faux positifs \emph{candidats} supplémentaires côté LLM.
Ce dernier consomme 1\,394 secondes, 1,34 dollar comptabilisé et
1\,189\,300 tokens. Une autre paire S1 est annoncée à égalité, 66,7\,\%.
Ces nombres sont attribués au journal ; ils ne résultent pas d'un nouveau
calcul ni d'une nouvelle inspection exhaustive de leurs traces
\cite{bench-preflight}.

Un changement de l'instrument de bilan modifie ensuite l'empreinte des
ressources. Les essais antérieurs sont conservés hors du pilote gelé.
Dans la copie locale consultée, les vingt-quatre chemins de bras restent
nuls et les précontrôles \code{pending}. Aucun résultat de campagne admissible
n'y est donc enregistré. Un workflow réussi ou un prochain identifiant
d'essai ne comble pas ce manque ; le bilan doit suivre les preuves effectivement
référencées \cite{bench-pilote,bench-preflight}.

\subsection{Plan V2 : rendre les questions et les mesures explicites}

La proposition V2 révisée le 15 octobre décrit 44 fiches
\code{s_...}, 25 ensembles techniques, 216 configurations sélectionnées et
456 bras conditionnels par instance et répétition communes ; 24 configurations
incluent A0. Une fiche décrit une situation, un type son traitement, un bras
la politique, une instance un monde paramétré et une répétition un nouvel
essai après reset. Ces nombres ne sont ni des exécutions accomplies, ni des
familles indépendantes à additionner au corpus public \cite{bench-v2}.

Ces unités se comprennent sur \code{s_ssh_identifiants} : la fiche pose la
question des identifiants faibles ; deux configurations retiennent une posture
protégée et une posture vulnérable. Pour une même instance, c'est-à-dire les
mêmes paramètres de cible et d'objet témoin, chaque configuration est exécutée
avec D1 et A1 dans l'ordre déclaré. Cela donne quatre bras par instance et
répétition communes. Une nouvelle répétition exige une remise à zéro et ne crée
pas une nouvelle famille. Les 456 bras comptent des traitements prévus avant
leur multiplication éventuelle par le nombre d'instances et de répétitions ;
ils ne représentent pas 456 réseaux ni 456 résultats acquis.

Six questions organisent la lecture : établir une preuve ; comprendre contexte
et droits ; interpréter un contrat nouveau ; suivre état et séquence ;
planifier un parcours contraint ; réviser après changement.
Elles ne définissent pas six scores cognitifs. Les cinq tâches du
tableau~\ref{tab:bench-v2} fixent des unités différentes.
Un cas protégé exige des contrôles réellement exécutés et un témoin légitime
fonctionnel ; l'abstention peut être correcte lorsque les observations sont
insuffisantes. Les comparaisons hors ligne emploient un interpréteur à règles
et un interpréteur LLM sur le même dossier figé \cite{bench-v2,bench-v2-plan}.

\begin{table}[htbp]
\centering\small
\caption{Cinq tâches V2 proposées ; les mesures sont à calibrer puis geler
avant l'évaluation.}
\label{tab:bench-v2}
\begin{tabularx}{\linewidth}{@{}>{\raggedright\arraybackslash}p{2.3cm}>{\raggedright\arraybackslash}p{4.2cm}Y@{}}
\toprule
Tâche & Mesure primaire proposée & Exemple et précaution\\
\midrule
M-T : ciblée & Proportion d'essais avec objectif évalué et preuve valide au budget &
SSH faible ou lecture MQTT privée ; contrôler aussi la variante protégée.\\
M-G : audit global & Nombre de défauts confirmés par essai au budget fixé, zéros inclus &
Sauvegardes privées ; déclarer erreurs, couverture et coûts.\\
M-I : interprétation & Proportion de conclusions conformes au dossier, abstention incluse &
Preuves contradictoires ; aucune exploration réseau n'est mesurée.\\
M-E : état & Proportion de diagnostics corrects de l'état observable &
Interruption OTA ; ne pas demander de deviner un état indiscernable.\\
M-C : couverture & Proportion de diagnostics étayés dans le graphe autorisé fini &
Cible inaccessible ; distinguer accès, évaluation et impossibilité de conclure.\\
\bottomrule
\end{tabularx}
\end{table}

L'ambition réseau reste à qualifier : 35 fiches sont uniquement locales et
36 uniquement à carte complète.

Local signifie que la cible est directement accessible depuis le point
d'entrée autorisé, sans relais réseau à organiser. Une carte complète signifie
que l'architecture utile est fournie dès le départ. Ces situations peuvent
conserver des difficultés de preuve, de droits ou de séquence, sans mesurer
la découverte d'une architecture inconnue. Les deux nombres décrivent des
axes différents : une fiche peut être à la fois locale et à carte complète.
Ils ne doivent donc pas être additionnés. Les limites d'accès, resets, coûts, chemins
témoins et oracles restent à attester. Le plan ne constitue donc pas une
nouvelle campagne déployée. Le programme de soumission doit être aligné sur
ces tâches : un critère générique en vrais positifs ne couvre pas automatiquement
un diagnostic d'état ou d'interprétation. Une égalité ou une dégradation du
LLM demeure une issue recevable \cite{bench-v2,bench-objectif}.

\subsection{Premier incrément SSH et conditions de passage à la mesure}

L'exemple SSH cherche à établir une lecture privée obtenue avec un couple
utilisateur--mot de passe faible, au-delà d'un port acceptant une connexion.
La fixture est l'environnement de test préparé ; le canari est un contenu
témoin propre à l'essai ; l'oracle est le dossier privé indiquant la posture
attendue et l'empreinte de ce contenu. Un compte légitime vérifie que le serveur
et l'objet sont fonctionnels. La politique reçoit la cible, ce compte, le
dictionnaire public et le chemin autorisé, mais ni le couple gagnant éventuel
ni le contenu attendu. L'évaluateur peut ainsi contrôler l'observation sans
fournir sa réponse à la politique \cite{bench-ssh,bench-ssh-code}.

Le premier incrément V2 implémente un contrat public SSH, un adaptateur D1
réutilisant l'exécuteur et un évaluateur hors ligne. Il réalise un contrôle
légitime puis parcourt un dictionnaire public fini ; identité effective et
lecture du même objet doivent être observées ensemble. L'oracle privé contient
la posture attendue et l'empreinte d'un canari propre à l'essai. Pour le cas
protégé F0, tous les couples doivent être explicitement refusés avec un témoin
fonctionnel. Pour le cas vulnérable F1, l'accès faible et la lecture privée
doivent être prouvés. Ici F1 désigne la posture, distincte du score \(F_1\)
\cite{bench-ssh,bench-ssh-code}.

La sortie annonce toujours \code{mt_admissible: false} : c'est un verdict
probatoire, pas encore une mesure ciblée admise. La fixture réelle F0/F1,
le raccordement A1, la parité des budgets, le reset et le gel expérimental
restent à produire. La validation consignée le 15 octobre rapporte 32 tests
SSH réussis et 3\,885 tests réussis, un ignoré, dans la suite complète.
Les processus SSH y sont simulés. Ces résultats antérieurs sont cités ;
aucun test n'a été relancé pour ce bilan \cite{bench-ssh}.

La CI prévoit validation des compositions, syntaxe Ansible, tests et contrôle
de l'image isolée. Elle qualifie des contrats logiciels, sans prouver toutes
les propriétés des services déployés. La priorité est de fermer ce premier
contraste, puis MQTT lecture, autorisation API, reprise MQTT et cache
inter-identités. Avant toute revendication, il faut enfin réconcilier manifeste,
ressources, mesures et analyse statistique, archiver les preuves et revoir les
désaccords. Le passage aux scénarios suivants dépend d'une qualification
technique, sans exiger une victoire du LLM \cite{net-ci,bench-v2}.


\section{Avancement transversal et portefeuille de recherche}
\label{sec:synthese}
\subsection{Une progression à lire par livrables}
Le projet a produit une chaîne technique qui permet de décrire un réseau,
d'orchestrer des contrôles, de conserver des observations, d'appliquer un
évaluateur et de préparer l'apprentissage à partir des erreurs. Ces réalisations
sont reliées : la qualité des données d'entraînement dépend des résultats
d'audit ; la valeur d'un score dépend de l'environnement et des preuves ;
l'utilité de l'interface dépend de l'accès à ces éléments.

Le tableau~\ref{tab:synthese-maturite} rassemble les dimensions d'avancement.
Il distingue les réalisations présentes dans le dépôt des validations
encore nécessaires : la présence du code ne démontre ni robustesse générale
ni avantage expérimental.

\begin{table}[htbp]
\centering\small
\begin{tabularx}{\linewidth}{@{}>{\raggedright\arraybackslash}p{27mm}YY@{}}
\toprule
Dimension & Réalisation vérifiable dans le dépôt & Validation encore nécessaire\\
\midrule
Modélisation et moteur & Graphe, contexte par run, six phases, outils, profils, artefacts et suivi de consommation. & Robustesse sur des configurations et tailles variées, avec preuves de couverture conservées.\\
Laboratoire & Catalogue, topologies, packs, playbooks, précontrôles et réservation partagée. & Qualification des propriétés réellement déployées, remises à zéro et scénarios retenus.\\
Sécurité & Contrats d'attribution, validations de preuve, statuts et frontières d'exécution. & Revues de preuves et analyse des cas non couverts sans généraliser les garanties.\\
Application & Lancement, suivi, arrêt, historique, configuration et parcours de comparaison. & Évaluation des usages et du travail humain sur tâches définies.\\
Modèles & Préparation, entraînement paramétrable, service multi-adaptateurs et routage par rôle. & Artefacts d'entraînement à vérifier ; comparaison du modèle de base et des spécialisations sur données indépendantes.\\
Évaluation & Évaluateur commun, métriques, comparateur apparié et bilan de campagne. & Campagne complète gelée et analyse alignée avec les hypothèses préenregistrées.\\
\bottomrule
\end{tabularx}
\caption{Maturité qualitative des ensembles du projet. Les conditions de
validation indiquent ce qu'il reste à établir, sans score global d'achèvement.}
\label{tab:synthese-maturite}
\end{table}

Le rapport de jalon III du programme G6259 présente la méthodologie de test
M2.3 et la plateforme virtuelle parmi les contributions de suivi. La question
scientifique d'un avantage du LLM constitue une preuve supplémentaire à
construire. L'achèvement déclaré d'un livrable méthodologique peut donc
coexister avec des campagnes expérimentales et des validations spécialisées
encore en cours. Le présent bilan ne réévalue pas les jalons contractuels des
partenaires ; il précise ce que les éléments locaux permettent d'examiner
\cite{intro-milestone}.

\subsection{Les recherches qui accompagnent le développement}
Plusieurs études soutiennent les évolutions du logiciel. Elles ont leur
question et leur statut propres ; leurs propositions ne sont pas toutes des
fonctionnalités livrées. L'étude sur le contexte cherche à conserver les
observations, les contradictions et les décisions utiles dans les limites de
requête. Les reproductions synthétiques du 29 septembre ont mis en évidence
des pertes possibles ; leur fréquence et leur effet sur les performances
réelles restent à mesurer \cite{syn-context}.

L'étude sur les méthodes d'audit examine une organisation plus explicite des
propriétés, de la portée et des obligations de preuve. Elle propose des
comparaisons hors ligne entre représentations des mêmes observations et une
relecture sémantique bornée. Ces propositions peuvent améliorer la façon de
discuter un résultat sans démontrer par elles-mêmes une meilleure détection
réseau \cite{syn-methodes}.

Le diagnostic d'organisation du code du 3 octobre a formulé des lots de
réorganisation et des critères d'acceptation. Son index indique ensuite une
mise en œuvre partielle, notamment des contextes de phase isolés. Les mesures
et descriptions de l'ancienne copie restent historiques ; l'architecture
actuelle se lit dans le code et les contrats courants. Cette distinction
permet de conserver le raisonnement du diagnostic sans maintenir deux
descriptions concurrentes du système \cite{syn-organisation}.

Le recentrage du corpus développe une carte de situations où l'interprétation,
la planification et la révision peuvent avoir des effets différents. Il affine
les unités de mesure au lieu de traiter tous les scénarios comme une difficulté
unique. Ce travail de conception et le premier incrément SSH constituent des
livrables ; la réalisation complète des environnements et des comparaisons
reste à conduire : serveurs réels, témoins vulnérables et protégés, remise
en état, raccordement de la politique LLM et campagne appariée \cite{bench-v2}.

Une spécialisation vers l'Internet des objets médicaux (IoMT) a également été
formulée comme sujet de travail. Le dossier est explicitement au statut de
proposition non démarrée. Il envisage des réseaux médicaux simulés et une
évaluation par rapport à la version généraliste. Il ne démontre ni support
clinique actuel ni sûreté sur dispositif médical réel. Cette piste élargit les
possibilités de recherche ; elle ne doit pas être comptée parmi les capacités
déjà qualifiées \cite{syn-iomt}.

\section{Dépendances, limites et priorités de poursuite}
\label{sec:priorites}
\subsection{Le chemin critique relie environnement, preuve et décision}
Le principal enjeu transversal est l'attribution. Une dégradation peut provenir
du laboratoire, de l'inventaire, du scanner, d'une décision, d'un budget ou de
l'assemblage du résultat. Pour corriger le bon composant et comparer deux
politiques, il faut conserver l'endroit où l'exécution s'est arrêtée et les
observations réellement disponibles.

La figure~\ref{fig:syn-dependances} représente la dépendance minimale d'une
conclusion comparative. Chaque étape possède ses livrables et son contrôle.
Une augmentation du nombre de scénarios ou de modèles ne compense pas un
défaut d'état initial, une preuve mal attribuée ou un résultat manquant.

\begin{figure}[htbp]
\centering
\begin{tikzpicture}
\node[box,text width=29mm,minimum height=17mm] (a) {Environnement\\et état initial qualifiés};
\node[box,text width=29mm,minimum height=17mm,right=8mm of a] (b) {Exécution\\et observations conservées};
\node[focal,text width=29mm,minimum height=17mm,right=8mm of b] (c) {Preuves\\et résultats évaluables};
\node[box,text width=29mm,minimum height=17mm,right=8mm of c] (d) {Comparaison\\et portée discutées};
\draw[flow] (a) -- (b);
\draw[flow] (b) -- (c);
\draw[flow] (c) -- (d);
\end{tikzpicture}
\caption{La validité d'une comparaison repose sur une chaîne de dépendances.
Les performances d'une politique sont interprétables lorsque les conditions
qui rendent sa mesure valide sont établies.}
\label{fig:syn-dependances}
\end{figure}

Les contraintes de ressources sont également transversales. Une requête longue
peut réduire le temps restant pour la vérification ; une projection trop
agressive peut masquer une observation ; une reprise peut augmenter la
consommation tout en permettant un livrable valide. Le compromis entre coût,
latence et couverture doit être mesuré sur des configurations identifiées,
avec les tentatives infructueuses conservées.

\subsection{Un programme de travail par critères de sortie}
Les priorités du tableau~\ref{tab:syn-suite} sont proposées à partir du bilan.
Elles concernent autant la plateforme que les expériences. Un critère de
sortie permet d'annoncer une avancée concrète et d'éviter une progression
estimée à partir du seul volume de développement.

\begin{table}[htbp]
\centering\small
\begin{tabularx}{\linewidth}{@{}>{\raggedright\arraybackslash}p{24mm}YY@{}}
\toprule
Priorité & Travail à mener & Critère de sortie proposé\\
\midrule
État reproductible & Réconcilier version, ressources, inventaire, état initial et manifeste. & Une exécution peut être reliée à une configuration et à des précontrôles archivés.\\
Qualification & Vérifier les propriétés des instances retenues et leur remise à zéro. & Témoins vulnérables/protégés et effets attendus réellement contrôlés.\\
Preuves et revue & Compléter les dossiers et examiner les cas rejetés ou indéterminés. & Conclusions attribuées aux bonnes propriétés, avec réserves et contre-preuves visibles.\\
Campagne & Exécuter les essais déclarés et publier leur bilan complet. & Bras présents, absents et échoués comptés ; analyse alignée au protocole gelé.\\
Spécialisation & Comparer modèle de base, adaptateurs et routage sur tâches communes. & Qualité technique, fiabilité, latence et coûts attribuables à la configuration étudiée.\\
Usage et partage & Qualifier les parcours et préparer un extrait reproductible. & Procédure exécutable, limites de portabilité explicites et retour d'utilisateurs sur tâches définies.\\
\bottomrule
\end{tabularx}
\caption{Feuille de route proposée pour consolider conjointement logiciel,
expérimentation et usages. Chaque critère demande des configurations
et des observations archivées.}
\label{tab:syn-suite}
\end{table}

La calibration du pilote, l'exécution de ses douze paires et la qualification du
premier scénario SSH V2 sont deux travaux liés mais distincts. Une campagne
doit conserver sa configuration pendant l'exécution ; les nouveaux scénarios,
outils et modèles doivent être introduits dans des lots identifiés, avec leur
propre gel. La préparation des recherches futures peut continuer à condition
de préserver cette attribution \cite{bench-pilote,bench-ssh}.

\subsection{Publication et diffusion}
Le dossier d'objectifs retient une décision interne le 16 novembre 2026, puis
une soumission visée à AsiaCCS 2027 le 11 décembre 2026. Ces dates sont les
jalons consignés du projet ; elles ne sont pas une garantie de résultats ni
une vérification nouvelle du calendrier de la conférence. Les contributions
envisagées portent sur le benchmark, les conditions d'écart et le coût de
l'avantage éventuel. Un résultat nul mais mesuré avec une précision suffisante
peut conduire à repositionner la contribution \cite{bench-objectif}.

La diffusion exige aussi un artefact compréhensible hors de la machine de
développement. Le laboratoire Proxmox, les fournisseurs locaux et certaines
ressources matérielles limitent la portabilité. Un extrait reproductible
doit préciser ce qu'il reproduit, ce qu'il simule et ce qui nécessite
l'infrastructure de référence. Les sources, traces, configurations et modes
d'emploi sont donc des livrables de recherche au même titre que les tableaux
de résultats.

\section{Conclusion générale}
LANCE rassemble plusieurs réalisations qui répondent à une même difficulté :
transformer des décisions de modèles et des contrôles réseau en résultats
examinables. Le moteur en six phases organise le travail ; le graphe décrit
les relations et soutient la priorisation ; le laboratoire produit des
situations contrôlées ; les contrats de preuve délimitent les conclusions ;
l'application rend l'exécution consultable ; l'entraînement prépare des
modèles adaptés aux rôles ; le benchmark fournit le cadre de comparaison.

L'avancement dépasse donc la seule rédaction d'un protocole. Des composants
logiciels, des scénarios, des essais historiques, des diagnostics et des outils
de qualification existent. Leur portée reste différente : des tests simulés
ne remplacent pas une qualification de laboratoire, une représentation de
chemin potentiel ne prouve pas une intrusion et une observation isolée ne
permet pas une conclusion générale sur le LLM.

La suite consiste à relier ces acquis dans des campagnes complètes, à qualifier
la spécialisation et à rendre les résultats reproductibles. L'objectif est de
pouvoir expliquer quelles propriétés ont été examinées, quelles preuves ont
été obtenues, où l'exécution a rencontré ses limites et dans quelles conditions
les modèles apportent un bénéfice utile. Le projet dispose des principaux
instruments pour poser ces questions ; leur réponse demande encore des
mesures contrôlées et une discussion de leur portée.

\begin{thebibliography}{999}
\addcontentsline{toc}{section}{Références}
\small\interlinepenalty=10000
\bibitem{intro-milestone} Programme NATO SPS G6259. \source{G6259_Milestone3_UMONS_FINAL.docx}{Rapport local de jalon III du 1er septembre 2026 ; contexte institutionnel déclaré.}
\bibitem{intro-readme} Projet LANCE. \source{README.md}{Objectifs et périmètre actuels, consultation du 15 octobre 2026.}
\bibitem{intro-protocol} Projet LANCE. \source{research/2026-09-27-agent-vs-automation/protocol.md}{Protocole proposé et comparaison des politiques ; copie locale au 15 octobre 2026.}
\bibitem{intro-mission} Projet LANCE. \source{metric/squelette-latex/parts/partie1.tex}{Contexte et mission de spécialisation HMoE ; document local consulté pour l'historique des objectifs.}
\bibitem{intro-q4} Équipe du projet. \source{report/q4-2025.tex}{Rapport historique Q4 2025 ; architecture de l'époque, sans reprise des performances comme résultats actuels.}
\bibitem{intro-q1} Équipe du projet. \source{report/q1-2026.tex}{Rapport historique Q1 2026 ; trajectoire architecturale déclarée, distincte de la version courante.}
\bibitem{intro-phases} Projet LANCE. \source{src/agent/phases/README.md}{Organisation courante des phases et contextes ; lecture du 15 octobre 2026.}
\bibitem{intro-research} Projet LANCE. \source{research/README.md}{Index et statut des études ; consultation du 15 octobre 2026.}
\bibitem{arch-pipeline} Projet LANCE. \source{src/agent/pipeline.py}{Cycle de vie, prérequis, choix de modèle, finalisation et mémoire ; lecture statique du 15 octobre 2026, sans exécution.}
\bibitem{arch-runner} Projet LANCE. \source{src/agent/core/runner.py}{Boucle de phase, résolution des outils et transactions de sauvegarde ; code local inspecté le 15 octobre 2026.}
\bibitem{arch-graphe} Projet LANCE. \source{src/graph_backend.py}{Interface GraphBackend et implémentation NetworkXBackend ; sémantique des requêtes examinée le 15 octobre 2026.}
\bibitem{arch-db} Projet LANCE. \source{src/db/database.py}{Schéma SQLite des fournisseurs, modèles, runs, scores et consommations ; lecture statique du 15 octobre 2026.}
\bibitem{arch-contexte} Projet LANCE. \source{src/agent/phases/context.py}{PhaseContext, copie des variables et services explicites ; code local au 15 octobre 2026.}
\bibitem{arch-registre} Projet LANCE. \source{src/agent/registry.py}{Responsabilités, prérequis, groupes d’outils et noms de livrables des phases ; confrontés aux modules courants le 15 octobre 2026.}
\bibitem{arch-chargeur} Projet LANCE. \source{src/loader.py}{Chargement YAML des objets d’infrastructure et choix de NetworkX par défaut ; lecture du 15 octobre 2026.}
\bibitem{arch-risque} Projet LANCE. \source{src/risk_scorer.py}{Politique heuristique directed-v2, exposition, centralité et combinaison des composantes ; code local du 15 octobre 2026.}
\bibitem{arch-chemins} Projet LANCE. \source{src/attack_path.py}{Pondération heuristique et chemins dirigés calculés par Dijkstra ; mécanisme présent, distinct d’une observation de pivot.}
\bibitem{arch-cve} Projet LANCE. \source{src/cve_lookup.py}{Recherche NVD et classification produit/version/CPE ; code inspecté le 15 octobre 2026, aucune interrogation distante pour ce rapport.}
\bibitem{arch-pondere} Projet LANCE. \source{src/agent/tools/disbalance_engine.py}{Graphe pondéré, coûts et comparaison des distances ; modèle de priorisation, non résultat expérimental.}
\bibitem{arch-outils-graphe} Projet LANCE. \source{src/agent/tools/graph_tools.py}{Parcours infrastructure, scénario public et découverte ; outils de consultation et absence de précalcul risque/chemins en scénario public.}
\bibitem{arch-graphe-contexte} Projet LANCE. \source{src/agent/core/graph_context.py}{État de graphe par run, liaison des workers et session autonome de compatibilité ; lecture du 15 octobre 2026.}
\bibitem{arch-recon} Projet LANCE. \source{src/agent/phases/recon/run.py}{Projection des observations et contrôles de couverture de reconnaissance ; code local au 15 octobre 2026.}
\bibitem{arch-analyse} Projet LANCE. \source{src/agent/phases/analysis/run.py}{Coordination scanner, analyses par équipement et agrégation ; lecture statique du 15 octobre 2026.}
\bibitem{arch-agregation} Projet LANCE. \source{src/agent/phases/analysis/aggregation.py}{Capture du contexte et agrégation conservant le registre brut ; code local au 15 octobre 2026.}
\bibitem{arch-verification} Projet LANCE. \source{src/agent/phases/verification/run.py}{Workers de vérification et reconstruction des statuts depuis les traces ; lecture du 15 octobre 2026.}
\bibitem{arch-intrusion} Projet LANCE. \source{src/agent/phases/intrusion/run.py}{Contexte et clôture de la phase 5, garde contre l’absence de tentative ; code local inspecté le 15 octobre 2026.}
\bibitem{arch-rapport} Projet LANCE. \source{docs/architecture/report-writing.md}{Rédaction par fiches et assemblage factuel ; limites conservées, anciens délais non utilisés comme contrat courant.}
\bibitem{arch-rapport-code} Projet LANCE. \source{src/agent/phases/report/run.py}{Rédaction séquentielle, réutilisation des fiches et délai adapté à leur nombre ; code du 15 octobre 2026.}
\bibitem{arch-prompts} Projet LANCE. \source{src/agent/prompt_manager.py}{Chargement des prompts, inclusions et interpolation ; lecture statique du 15 octobre 2026.}
\bibitem{arch-outils} Projet LANCE. \source{src/agent/tools/tool_loader.py}{Schémas YAML, adaptateurs, suggestions matérielles et filtrage de disponibilité ; code courant, sans vérification d’équipement réel.}
\bibitem{arch-provider} Projet LANCE. \source{src/agent/provider.py}{Adaptateur de fournisseurs compatibles OpenAI et boucle synchrone d’appels d’outils ; lecture du 15 octobre 2026.}
\bibitem{arch-executeur} Projet LANCE. \source{src/agent/core/executor.py}{Exécuteur instrumenté, budgets et politiques de connaissances ; code local au 15 octobre 2026.}
\bibitem{arch-profils} Projet LANCE. \source{src/agent/execution_profiles.py}{Profils compact/full, routage des outils et résolution auto au seuil de 32 milliards ; configuration inspectée le 15 octobre 2026.}
\bibitem{arch-recon-lots} Projet LANCE. \source{src/agent/phases/recon/subagents.py}{Balayage, lots de couverture et fusion depuis le journal ; conditions d’éligibilité examinées le 15 octobre 2026.}
\bibitem{arch-reprise-graphe} Projet LANCE. \source{docs/architecture/phase1-graph-recovery.md}{Reprise de rédaction après troncature en phase 1 ; mécanisme présent, sans nouveau rejeu expérimental.}
\bibitem{arch-reprise-blocs} Projet LANCE. \source{docs/architecture/phase3-block-recovery.md}{Reprise par blocs de services et conditions de promotion ; documentation consultée le 15 octobre 2026.}
\bibitem{arch-connaissances} Projet LANCE. \source{src/agent/tools/skill_tools.py}{Fiches, recherche sémantique, classification CVE et instantané du benchmark ; lecture du 15 octobre 2026.}
\bibitem{arch-memoire} Projet LANCE. \source{src/agent/knowledge/ingest.py}{Ingestion de méthodes et de constats ; exclusion des artefacts benchmark de la mémoire épisodique, lecture du 15 octobre 2026.}
\bibitem{arch-transport} Projet LANCE. \source{src/agent/core/provider_transport.py}{Reprises de transport et délais distincts des continuations modèle ; code inspecté le 15 octobre 2026.}
\bibitem{arch-admission} Projet LANCE. \source{src/agent/core/provider_concurrency.py}{Admission des requêtes par origine et libération du verrou ; code local du 15 octobre 2026.}
\bibitem{arch-cout} Projet LANCE. \source{src/agent/cost_tracker.py}{Mesures par phase, provenance des tarifs, durée cumulée et écoulée, budget de consommation observée ; lecture du 15 octobre 2026.}
\bibitem{arch-stockage} Projet LANCE. \source{src/agent/knowledge/store.py}{Client ChromaDB persistant, recherche et repli éphémère ; code inspecté le 15 octobre 2026, sans ouverture d'une base pour ce rapport.}
\bibitem{arch-dispatch} Projet LANCE. \source{src/agent/phases/registry.py}{Routage explicite des six modules de phase ; code local inspecté le 15 octobre 2026.}
\bibitem{net-scope} Projet LANCE. \source{AGENTS.md}{Cible nato/pve-nato, runner nato-master et exclusion du homelab personnel ; instructions locales lues le 15 octobre 2026.}
\bibitem{net-master} Projet LANCE. \source{benchmarks/ansible/playbooks/deploy_master.yml}{Provisionnement du contrôleur LXC, interfaces, filtrage, Tailscale et service applicatif ; lecture statique ciblée le 15 octobre 2026. Aucun état distant vérifié.}
\bibitem{net-deploy} Projet LANCE. \source{benchmarks/ansible/playbooks/03_deploy_scenario.yml}{Clonage, appartenance des charges, conflits, interfaces VLAN et préparation réseau ; lecture statique le 15 octobre 2026.}
\bibitem{net-access} Projet LANCE. \source{docs/guides/lab-access-tailscale.md}{Accès au laboratoire ; procédure documentée comme vérifiée le 7 octobre 2026, relue le 15 octobre sans nouvelle connexion.}
\bibitem{net-scenarios} Projet LANCE. \source{docs/benchmark/v1/scenarios.md}{Corpus, séparation des groupes, composition, exposition de la topologie et limites de déploiement ; état local consulté le 15 octobre 2026.}
\bibitem{net-attacks} Projet LANCE. \source{docs/benchmark/catalogue-attaques.md}{Guide de couverture et limites des simulations ; relecture documentaire datée du 15 septembre 2026, confrontée aux fichiers locaux le 15 octobre.}
\bibitem{net-topology} Projet LANCE. \source{benchmarks/topologies/true_multihop.yaml}{Objectif S20 à deux relais et segments isolés ; définition locale lue le 15 octobre 2026, sans essai réseau.}
\bibitem{net-verify} Projet LANCE. \source{benchmarks/ansible/playbooks/06_verify.yml}{Contrôles de propriétés, refus attendus, chemins réseau et consolidation du preflight ; lecture statique ciblée le 15 octobre 2026, sans exécution.}
\bibitem{net-ot} Projet LANCE. \source{benchmarks/packs/definitions/f15_ot_protocols.yaml}{Propriétés et contrôles des simulateurs OT et registre Modbus jetable ; définition locale consultée le 15 octobre 2026.}
\bibitem{net-catalog} Projet LANCE. \source{benchmarks/catalog.yaml}{Catalogue 3.2.0 : 29 scénarios publics et groupes ; lecture complète le 15 octobre 2026.}
\bibitem{net-manual} Projet LANCE. \source{docs/guides/manual-scenarios.md}{Composition manuelle, profils, adaptateurs d'exécution, prévisualisation et sémantique des chemins ; lecture ciblée le 15 octobre 2026.}
\bibitem{net-generated} Projet LANCE. \source{src/benchmark/scenario_deployment.py}{Adaptation des exports aux playbooks, surcouches et réservations d'identifiants ; lecture ciblée le 15 octobre 2026.}
\bibitem{net-lifecycle} Projet LANCE. \source{src/agent/core/lifecycle.py}{Cycle déploiement--injection--vérification, contrat preflight-v1 et nettoyage ; lecture ciblée le 15 octobre 2026.}
\bibitem{net-reset} Projet LANCE. \source{benchmarks/ansible/playbooks/08_reset_scenario.yml}{Redémarrage, réinjection et repeuplement ; lecture complète le 15 octobre 2026, sans reset exécuté.}
\bibitem{net-instances} Projet LANCE. \source{docs/guides/development-instances.md}{Trois emplacements applicatifs, réservation commune, activité déclarée, CI et limites ; état documentaire local relu le 15 octobre 2026.}
\bibitem{net-lock} Projet LANCE. \source{src/benchmark/lab_lock.py}{Réservation coopérative interprocessus du laboratoire ; lecture ciblée le 15 octobre 2026.}
\bibitem{net-ci} Projet LANCE. \source{.github/workflows/benchmark-integrity.yml}{Validation complète ou documentaire et conditions d'appel du déploiement ; état local lu le 15 octobre 2026.}
\bibitem{net-update} Projet LANCE. \source{.github/workflows/update-master.yml}{Déploiement du commit validé sur le runner nato-master ; état local lu le 15 octobre 2026. Le script scripts/deployment/deploy-instance.sh a également été consulté pour ses contrôles d'identité et d'emplacement.}
\bibitem{sec-identite} Projet LANCE. \source{src/agent/finding_identity.py}{Identité conservatrice des constats et séparation des identifiants de provenance ; lecture statique du 15 octobre 2026.}
\bibitem{sec-evaluateur} Projet LANCE. \source{src/benchmark/evaluator.py}{Chargement des traces et attribution des preuves, notamment \code{_tool_call_matches_finding} et \code{_assign_evidence_refs} ; lecture ciblée du 15 octobre 2026, sans calcul nouveau.}
\bibitem{sec-verification} Projet LANCE. \source{src/agent/phases/verification/contract.py}{Plans de vérification par famille et routage des outils ; lecture ciblée du 15 octobre 2026.}
\bibitem{sec-preuves} Projet LANCE. \source{src/agent/exploit_evidence.py}{Interprétation partagée des traces et synthèse des verdicts ; lecture ciblée du 15 octobre 2026.}
\bibitem{sec-ssh} Projet LANCE. \source{src/agent/evidence/ssh.py}{Politique finie \code{ssh-crypto-policy-v1}, preuves d'algorithmes observés ; lecture du 15 octobre 2026.}
\bibitem{sec-mqtt} Projet LANCE. \source{src/agent/evidence/mqtt.py}{Attribution des messages au topic, format structuré et compatibilité historique ; lecture ciblée du 15 octobre 2026.}
\bibitem{sec-evaluation} Projet LANCE. \source{docs/benchmark/v1/evaluation.md}{Contrat courant de preuve, métriques, résultats indéterminés et absence de provenance de pivot ; consultation du 15 octobre 2026.}
\bibitem{sec-scanner} Projet LANCE. \source{src/agent/scanner.py}{Extracteurs PKI et mises à jour signées, fonctions \code{_extract_pki_findings} et \code{_extract_ota_findings} ; lecture ciblée du 15 octobre 2026, sans sonde.}
\bibitem{sec-rapport} Projet LANCE. \source{src/agent/report_evidence.py}{Admission au rapport et distinction entre échec d'exécution et preuve négative ; lecture du 15 octobre 2026.}
\bibitem{sec-acces} Projet LANCE. \source{src/agent/evidence/capabilities.py}{Règles d'accès observé, distinctes du contact réseau ; lecture du 15 octobre 2026.}
\bibitem{sec-intrusion} Projet LANCE. \source{src/agent/phases/intrusion/observations.py}{Projection des accès corroborés, contrôles de journal et transitions indisponibles ; lecture du 15 octobre 2026.}
\bibitem{sec-gate} Projet LANCE. \source{docs/architecture/phase5-noop-gate.md}{Gate de tentative observable du profil complet, conditions et limites ; consultation du 15 octobre 2026.}
\bibitem{sec-scope} Projet LANCE. \source{src/agent/phases/intrusion/scope.py}{Contrôle applicatif des destinations directes et commandes distantes ; lecture du 15 octobre 2026.}
\bibitem{sec-auth} Projet LANCE. \source{src/api/auth.py}{Authentification administrative, exceptions de lancement et sessions révocables ; lecture du 15 octobre 2026.}
\bibitem{sec-session} Projet LANCE. \source{docs/guides/admin-session.md}{Session de huit heures, transport, réseau de confiance et limite à un processus ; consultation du 15 octobre 2026.}
\bibitem{sec-tests-frontieres} Projet LANCE. \source{tests/test_review_mqtt_ws_boundary.py}{Tests des contrôles de périmètre, arrêt et budget de la sonde WebSocket.} \source{tests/pipeline/test_verification_evidence.py}{Non-admission d'un upgrade WebSocket comme preuve d'accès MQTT.} \source{tests/test_intrusion_access_consistency.py}{Cohérence des règles d'accès entre outil, livrable compact et projection ; lectures ciblées du 15 octobre 2026, sans réexécution.}
\bibitem{sec-tests-identite} Projet LANCE. \source{tests/pipeline/test_pki_identity_evidence.py}{Contrôles positifs/négatifs et déduplication des équipements pour l'identité PKI ; lecture du 15 octobre 2026, sans réexécution.}
\bibitem{sec-tests-auth} Projet LANCE. \source{tests/test_review_mutation_auth.py}{Tests du refus des mutations administratives et de l'exception des lancements publics ; lecture du 15 octobre 2026, sans réexécution.}
\bibitem{sec-contrat} Projet LANCE. \source{src/benchmark/metric_contract.py}{Versions \code{strict-v3.16} et \code{evidence-v15}, empreintes des prompts et schémas d'outils ; lecture du 15 octobre 2026.}
\bibitem{app-main} Projet LANCE. \source{src/api/main.py}{Point d'entrée applicatif, routes et identité d'environnement ; état local consulté le 15 octobre 2026.}
\bibitem{app-runs} Projet LANCE. \source{src/api/routes/runs.py}{Historique, statuts, artefacts et projection des diagnostics ; code consulté le 15 octobre 2026, sans vérification du service distant.}
\bibitem{app-ui} Projet LANCE. \source{src/static/app.js}{Tableau de bord, graphe, synthèses et comparaison ; état local consulté le 15 octobre 2026.}
\bibitem{app-modes} Projet LANCE. \source{docs/architecture/run-modes.md}{Modes d'exécution et frontière avec l'application ; contrat courant consulté le 15 octobre 2026.}
\bibitem{app-start} Projet LANCE. \source{src/api/routes/pipeline.py}{Validation, lancement, batch, arrêt et suivi ; code consulté le 15 octobre 2026.}
\bibitem{app-providers} Projet LANCE. \source{src/api/routes/providers.py}{Contrat de configuration des fournisseurs ; code consulté le 15 octobre 2026.}
\bibitem{app-auth-guide} Projet LANCE. \source{docs/guides/provider-auth.md}{Authentification des mutations et fournisseurs exécutables ; guide courant consulté le 15 octobre 2026.}
\bibitem{app-artifacts} Projet LANCE. \source{docs/architecture/run-artifacts.md}{Isolation des livrables et séquencement des appels au modèle ; guide consulté le 15 octobre 2026. Les validations historiques citées dans ce guide ne sont pas de nouveaux essais du présent bilan.}
\bibitem{app-coordinator} Projet LANCE. \source{src/api/run_coordinator.py}{Coordination d'un run, état, arrêt et progression ; code consulté le 15 octobre 2026.}
\bibitem{app-events} Projet LANCE. \source{src/api/events.py}{Journal SSE borné et curseurs indépendants ; code consulté le 15 octobre 2026.}
\bibitem{app-activity} Projet LANCE. \source{src/api/routes/activity.py}{Vue d'activité en lecture seule des instances connues ; code consulté le 15 octobre 2026, sans interrogation des instances.}
\bibitem{app-batch-tests} Projet LANCE. \source{tests/test_review_batch_outcomes.py}{Tests des statuts, décomptes et coûts de batch ; tests lus, non exécutés pour cette contribution.}
\bibitem{app-start-tests} Projet LANCE. \source{tests/test_api_run_modes.py}{Tests des contrats de lancement et de coordination ; tests lus, non exécutés pour cette contribution.}
\bibitem{app-rules-tests} Projet LANCE. \source{tests/test_rules_audit_ui.py}{Tests de restitution des contrôles par règles ; tests lus, non exécutés pour cette contribution.}
\bibitem{app-intrusion-tests} Projet LANCE. \source{tests/test_intrusion_observations_ui.py}{Tests des projections d'accès et de l'absence de promotion des déclarations en preuves ; tests lus, non exécutés pour cette contribution.}
\bibitem{app-summary-tests} Projet LANCE. \source{tests/test_pipeline_summary_ui.py}{Tests des synthèses applicatives et des données indisponibles ; tests lus, non exécutés pour cette contribution.}
\bibitem{ml-memoire} Projet LANCE. \source{metric/squelette-latex/parts/partie3.tex}{Mémoire de stage, chapitre QLoRA et HMoE ; dernière révision Git a98b480 du 28 août 2026. Provenances, volumes et pertes historiques rapportés ; artefacts GPU non réinspectés.}
\bibitem{ml-config} Projet LANCE. \source{model_training/configs/qlora_qwen2_5_3b.yaml}{Profil 3B à quatre experts, hyperparamètres et feedback configurés ; copie locale lue le 15 octobre 2026, sans attestation d'exécution.}
\bibitem{ml-train} Projet LANCE. \source{model_training/train_qlora_3b.py}{Chaîne SFT, séparation de validation, injection du feedback, chargement quantifié, sélection du checkpoint ; code lu le 15 octobre 2026, sans entraînement.}
\bibitem{ml-template} Projet LANCE. \source{model_training/templates/qwen2_5_tools_assistant_mask.jinja}{Gabarit conversationnel Qwen et masques de génération des messages assistant ; lu le 15 octobre 2026.}
\bibitem{ml-prepare} Projet LANCE. \source{model_training/prepare_3b_datasets.py}{Fenêtrage, catalogue d'outils, compactage et chaîne atomique de phase 5 ; code lu le 15 octobre 2026.}
\bibitem{ml-preflight} Projet LANCE. \source{model_training/preflight_3b.py}{Contrôles de dépendances, données, CUDA et adaptateurs ; lecture du 15 octobre 2026, sans exécution GPU.}
\bibitem{ml-server} Projet LANCE. \source{src/agent/moe_server.py}{Service compatible OpenAI, routage, verrou de génération et budgets de contexte ; code lu le 15 octobre 2026, sans constat d'état distant.}
\bibitem{ml-service} Projet LANCE. \source{docs/training/hmoe-openwebui.md}{Procédure de service HMoE et d'accès privé ; consultée le 15 octobre 2026, sans connexion à une machine externe.}
\bibitem{ml-workspaces} Projet LANCE. \source{docs/training/workspaces.md}{Séparation du workspace canonique et du workspace d'entraînement, promotion et provenance ; consulté le 15 octobre 2026.}
\bibitem{ml-sync} Projet LANCE. \source{model_training/workspace_sync.json}{Liste blanche et limites de collecte des rapports et adaptateurs ; copie locale lue le 15 octobre 2026.}
\bibitem{ml-learning} Projet LANCE. \source{docs/training/learning-loop.md}{Minage des erreurs, revue humaine, attribution aux experts et exclusion des tests ; référence consultée le 15 octobre 2026.}
\bibitem{ml-sft} Projet LANCE. \source{src/learning/sft_feedback.py}{Conversion des corrections acceptées, contrôles de phase et d'origine ; code lu le 15 octobre 2026.}
\bibitem{ml-import} Projet LANCE. \source{src/learning/remote_runs.py}{Import contrôlé d'artefacts et manifeste de provenance ; code consulté le 15 octobre 2026, aucun import exécuté.}
\bibitem{ml-readme} Projet LANCE. \source{model_training/README.md}{Périmètre des outils et limites des tests sans entraînement GPU ; consulté le 15 octobre 2026.}
\bibitem{bench-composition} Projet LANCE. \source{benchmarks/scenarios/dev/S15.yaml}{Composition S15, lue avec la topologie, la vérité terrain scenario\_15.yaml et le pack f11\_api\_authorization.yaml ; descriptions statiques, sans nouvelle vérification du laboratoire.}
\bibitem{bench-labels} Projet LANCE. \source{docs/benchmark/derived-index.md}{Origine et limites des axes calculés depuis les vérités terrain ; consultation le 15 octobre 2026.}
\bibitem{bench-couverture} Projet LANCE. \source{objectifs/asiaccs-2027/audit-couverture.md}{Audit statique daté du 12 octobre 2026 ; 288 vulnérabilités attendues et proposition de pilote, sans run.}
\bibitem{bench-funnel} Projet LANCE. \source{src/benchmark/funnel.py}{Entonnoir candidats/retenus/confirmés et disponibilité des populations ; lecture statique du code local.}
\bibitem{bench-cli} Projet LANCE. \source{docs/benchmark/agent-vs-automation-cli.md}{Périmètre D1/A1 des phases 3--4, outils et calibration SSH ; la mention ancienne d'une inférence absente est confrontée au code courant.}
\bibitem{bench-policy-code} Projet LANCE. \source{src/agent/audit_experiment.py}{Périmètre des expériences sur inventaire et version shared-scanner-bounded-verification-v5 ; lecture statique du code local.}
\bibitem{bench-rules} Projet LANCE. \source{src/agent/phases/verification/rules.py}{Règles de reprise bornée, observations et validation de preuves communes ; lecture statique au 15 octobre 2026.}
\bibitem{bench-campagne} Projet LANCE. \source{docs/benchmark/agent-vs-automation-campaign.md}{Manifestes, issues opérationnelles, ressources, archivage et revue indépendante ; consultation locale au 15 octobre 2026.}
\bibitem{bench-comparateur} Projet LANCE. \source{src/benchmark/compare_policies.py}{Critères d'admissibilité d'une paire et export des motifs ; aucun comparateur exécuté pour ce rapport.}
\bibitem{bench-campaign-code} Projet LANCE. \source{src/benchmark/policy_campaign.py}{Bilan opérationnel, regroupement des configurations et appel de l'inférence multi-familles ; lecture statique du 15 octobre 2026.}
\bibitem{bench-inference} Projet LANCE. \source{src/benchmark/inference.py}{Bootstrap par famille, intervalle à 95 \% et Holm entre métriques ; implémentation présente, sans application à une campagne admissible dans ce bilan.}
\bibitem{bench-objectif} Projet LANCE. \source{objectifs/asiaccs-2027/plan-soumission.md}{Plan de soumission retenu le 12 octobre 2026 et critères en vrais positifs confirmés ; objectif à aligner aux tâches V2, sans résultat revendiqué.}
\bibitem{bench-qwen} Projet LANCE. \source{research/2026-10-05-benchmark-coverage/README.md}{Premier lot Qwen UMONS commencé le 4 octobre ; données gelées le 5 octobre 2026 à 11:32, rapport révisé à 11:48 ; valeurs historiques transcrites sans recalcul.}
\bibitem{bench-run1} Projet LANCE. \source{research/2026-10-05-corrections-lance/run-01.tex}{Rapport Run1 révisé le 5 octobre 2026 à 23:38 ; scores et statuts historiques, réserves sémantiques préservées.}
\bibitem{bench-run2} Projet LANCE. \source{research/2026-10-05-corrections-lance/run-02.tex}{Campagne corrective distincte, révision du 6 octobre 2026 à 00:49 ; lecture de métadonnées sélectionnées.}
\bibitem{bench-diagnostic} Projet LANCE. \source{research/2026-10-05-corrections-lance/diagnostic-pipeline.tex}{Diagnostic hors ligne révisé le 6 octobre 2026 à 10:34, inventaire de 09:29 ; 2 done, 8 partial, 2 failed ; sans certification exhaustive des preuves.}
\bibitem{bench-pilote} Projet LANCE. \source{benchmarks/experiments/policy-comparison/pilot-2026-10.json}{Manifeste réel lu directement le 15 octobre 2026 : 12 paires prévues, 24 run\_dir nuls et précontrôles pending ; configuration et budgets consignés.}
\bibitem{bench-preflight} Projet LANCE. \source{benchmarks/experiments/policy-comparison/pilot-2026-10-preflight.md}{Journal des préparatifs et fumées des 13--14 octobre 2026, calibration et changement de gel ; nombres attribués au journal, sans nouveau calcul des traces.}
\bibitem{bench-v2} Projet LANCE. \source{research/2026-10-12-recentrage-corpus/README.md}{Plan V2 révisé le 15 octobre 2026 à 14:49, proposition sans résultat expérimental nouveau : 44 fiches, 25 ensembles, 216 configurations et 456 bras conditionnels.}
\bibitem{bench-v2-plan} Projet LANCE. \source{research/2026-10-12-recentrage-corpus/coverage-plan.json}{Plan structuré proposé, contrats, tâches et comparateurs ; aucune exécution déduite des nombres de configurations.}
\bibitem{bench-ssh} Projet LANCE. \source{benchmarks/experiments/ssh-reference/README.md}{Premier incrément SSH et qualification logicielle consignée le 15 octobre 2026 ; SSH simulé, fixture réelle et A1 restant à réaliser.}
\bibitem{bench-ssh-code} Projet LANCE. \source{benchmarks/experiments/ssh_reference.py}{Contrat public, adaptateur D1 et évaluation probatoire SSH ; lecture du code courant, sortie non admise comme mesure M-T.}
\bibitem{syn-context} Projet LANCE. \source{research/2026-09-29-context-scalability/README.md}{Étude active, diagnostic et reproductions du 29 septembre 2026 ; impact réel non mesuré.}
\bibitem{syn-methodes} Projet LANCE. \source{research/2026-10-07-methodes-audit/README.md}{Exploration des méthodes et protocole proposé, révision du 7 octobre 2026.}
\bibitem{syn-organisation} Projet LANCE. \source{research/2026-10-03-organisation-code/README.md}{Diagnostic historique et index mentionnant une mise en œuvre partielle ultérieure.}
\bibitem{syn-iomt} Projet LANCE. \source{research/2026-09-29-iomt/README.md}{Proposition IoMT non démarrée ; date de classement distincte d'une date de début connue.}
\end{thebibliography}

\appendix
\section{Repères de vocabulaire}
\label{sec:glossaire}
Ce tableau rassemble les principaux termes employés dans les parties
spécialisées. Une confirmation désigne une déclaration soutenue par les
observations admises au contrat applicable. Elle reste limitée à la propriété,
à la cible et aux conditions effectivement contrôlées.
\small
\begin{longtable}{@{}>{\raggedright\arraybackslash}p{27mm}>{\raggedright\arraybackslash}p{128mm}@{}}
\toprule Terme & Sens dans le rapport\\\midrule\endfirsthead
\toprule Terme & Sens dans le rapport (suite)\\\midrule\endhead
\bottomrule\endfoot
API & Interface de programmation reliant le tableau de bord, le contrôleur et les services.\\
Benchmark & Banc d'essai avec configurations et critères de mesure identifiés.\\
Bras & Exécution d'une politique ou d'un traitement dans une comparaison.\\
CIDR & Notation d'une plage d'adresses réseau utilisée pour déclarer le périmètre.\\
CPE / CVE & Identifiant de produit ou de plateforme / identifiant d'une vulnérabilité publiée.\\
CVSS & Cotation de sévérité ; ne mesure pas la probabilité d'une compromission.\\
D1 / A1 / A0 & Automatisation conditionnelle sans LLM / politique LLM révisable / plan initial LLM figé. A0 est proposé ; D1/A1 dispose d'un incrément logiciel.\\
Ground truth & Vérité terrain : propriétés attendues réservées à l'évaluation.\\
HMoE & Combinaison de rôles spécialisés via adaptateurs, routés par phase au niveau de la requête.\\

Instance & Selon le contexte, processus applicatif ou monde expérimental paramétré. Dans une paire de benchmark, les politiques reçoivent les mêmes paramètres de monde.\\
IoT / OT / IoMT & Objets connectés / systèmes opérationnels industriels / objets connectés médicaux.\\
LLM & Modèle de langage de grande taille, employé pour certaines décisions et explications.\\
LoRA / QLoRA & Adaptation par paramètres supplémentaires de faible rang / entraînement de ces adaptateurs avec base quantifiée.\\
MQTT / topic & Protocole de messagerie publish/subscribe / canal de messages.\\
NVD & Base de vulnérabilités utilisée pour l'enrichissement, ou via un instantané en benchmark.\\
OTA & Mise à jour à distance ; le verdict reste limité aux propriétés contrôlées.\\
Pack / posture & Ensemble de propriétés composables / profil de configuration d'une instance.\\
Pivot & Action depuis un accès acquis vers une autre cible, avec causalité de la transition réseau à démontrer.\\
Preflight / reset & Vérifications avant audit / procédure de remise en état entre essais.\\
Run & Exécution identifiée, avec son dossier, ses contextes, ses observations et son état terminal.\\
SSE & Flux d'événements du serveur vers les lecteurs pour le suivi d'exécution.\\
SSH / TLS & Connexion distante chiffrée / protection du transport et authentification selon configuration.\\
VLAN & Séparation de domaines réseau ; routage et filtrage demandent des règles distinctes.\\
VP / FP / FN & Vrai positif / faux positif / faux négatif, définis selon la population et la vérité terrain.\\
\end{longtable}
\normalsize

\end{document}
